\documentclass[11pt]{article}% Shared preamble for the rigor documents (Lamport-style structured proofs).  \input this file after \documentclass.
\usepackage[T1]{fontenc}\usepackage{lmodern}\usepackage{microtype}
\usepackage[margin=1in]{geometry}
\usepackage{amsmath,amssymb,amsthm,mathtools}
\usepackage{xcolor}\usepackage{enumitem}\usepackage{booktabs}\usepackage{graphicx}\usepackage{hyperref}
\usepackage[most]{tcolorbox}
\definecolor{navy}{HTML}{17365D}\definecolor{teal}{HTML}{117A8B}\definecolor{warn}{HTML}{A33A2B}\definecolor{okgreen}{HTML}{2E7D4F}
\hypersetup{colorlinks=true,linkcolor=navy,citecolor=teal,urlcolor=teal}
\theoremstyle{plain}
\newtheorem{theorem}{Theorem}[section]\newtheorem{lemma}[theorem]{Lemma}\newtheorem{proposition}[theorem]{Proposition}\newtheorem{corollary}[theorem]{Corollary}
\theoremstyle{definition}\newtheorem{definition}[theorem]{Definition}\newtheorem{remark}[theorem]{Remark}\newtheorem{claim}[theorem]{Claim}\newtheorem{issue}[theorem]{Issue}
% ---------- Lamport structured proofs ----------
% \lstep{level}{number}{assertion}   -> indented "<level>number. assertion"
% \lproof                             -> "PROOF:" line (start of the sub-proof of the preceding step)
% \lby{justification}                 -> "BY ..." leaf justification for the preceding step
% \lqed{level}{number}                -> QED step of a sub-proof
% keywords: \lassume \lprove \lsuffices \lcase \ldefine \lpick \lnew
\newlength{\lindent}\setlength{\lindent}{1.7em}
\newcommand{\lstep}[3]{\par\smallskip\noindent\hspace*{\dimexpr\numexpr#1-1\relax\lindent\relax}{\color{navy}$\langle#1\rangle#2$.}\ #3}
\newcommand{\lproof}{\par\noindent\hspace*{\lindent}{\scshape Proof:}}
\newcommand{\lby}[1]{\par\noindent\hspace*{\lindent}{\scshape By}\ #1.}
\newcommand{\lqed}[2]{\par\smallskip\noindent\hspace*{\dimexpr\numexpr#1-1\relax\lindent\relax}{\color{navy}$\langle#1\rangle#2$.}\ {\scshape Q.E.D.}}
\newcommand{\lassume}{{\scshape Assume}\ }\newcommand{\lprove}{{\scshape Prove}\ }\newcommand{\lsuffices}{{\scshape Suffices}\ }
\newcommand{\lcase}{{\scshape Case}\ }\newcommand{\ldefine}{{\scshape Define}\ }\newcommand{\lpick}{{\scshape Pick}\ }\newcommand{\lnew}{{\scshape New}\ }
% ---------- verdict boxes ----------
\newtcolorbox{verdictok}{colback=okgreen!8,colframe=okgreen,title={\bfseries Verdict: verified},fonttitle=\small}
\newtcolorbox{verdictgap}{colback=orange!10,colframe=orange!80!black,title={\bfseries Verdict: gap / caveat},fonttitle=\small}
\newtcolorbox{verdictbreak}{colback=warn!10,colframe=warn,title={\bfseries Verdict: proof-breaking issue},fonttitle=\small}
\newtcolorbox{citedbox}[1]{colback=navy!5,colframe=navy!60,title={\bfseries Cited result: #1},fonttitle=\small,breakable}
\setlength{\parindent}{0pt}\setlength{\parskip}{4pt}\allowdisplaybreaks
\newcommand{\cA}{\mathcal A}\newcommand{\cF}{\mathcal F}\newcommand{\cM}{\mathfrak M}\newcommand{\cN}{\mathfrak N}

% extra calligraphic shortcuts (\providecommand: no clash if a document defines its own)
\providecommand{\cB}{\mathcal B}\providecommand{\cC}{\mathcal C}\providecommand{\cD}{\mathcal D}\providecommand{\cE}{\mathcal E}\providecommand{\cG}{\mathcal G}\providecommand{\cH}{\mathcal H}\providecommand{\cI}{\mathcal I}\providecommand{\cJ}{\mathcal J}\providecommand{\cK}{\mathcal K}\providecommand{\cL}{\mathcal L}\providecommand{\cO}{\mathcal O}\providecommand{\cP}{\mathcal P}\providecommand{\cQ}{\mathcal Q}\providecommand{\cR}{\mathcal R}\providecommand{\cS}{\mathcal S}\providecommand{\cT}{\mathcal T}\providecommand{\cU}{\mathcal U}\providecommand{\cV}{\mathcal V}\providecommand{\cW}{\mathcal W}\providecommand{\cX}{\mathcal X}\providecommand{\cY}{\mathcal Y}\providecommand{\cZ}{\mathcal Z}
\newcommand{\Fid}{\operatorname{F}}\newcommand{\Tr}{\operatorname{Tr}}\newcommand{\eps}{\varepsilon}\newcommand{\dd}{\,\mathrm d}


\newcommand{\Om}{\Omega}
\newcommand{\Hil}{\mathcal H}
\newcommand{\Msa}{\cM_{\mathrm{sa}}}
\newcommand{\Mpsa}{\cM'_{\mathrm{sa}}}
\newcommand{\HM}{\mathsf H_{\cM}}
\newcommand{\HMp}{\mathsf H_{\cM'}}
\newcommand{\Var}{\operatorname{Var}}
\newcommand{\dist}{\operatorname{dist}}
\newcommand{\gB}{g_{\mathrm B}}
\newcommand{\gKM}{g_{\mathrm{KM}}}
\newcommand{\ip}[2]{\langle #1,#2\rangle}
\newcommand{\norm}[1]{\lVert #1\rVert}
\newcommand{\U}{\mathcal U}
\newcommand{\one}{\mathbf 1}
\newcommand{\Ad}{\operatorname{Ad}}\newcommand{\ket}[1]{|#1\rangle}\newcommand{\bra}[1]{\langle#1|}

\title{\color{navy}\bfseries The type-III second-variation lemma\\[0.3em]
\large Note~7, Lemma~3.2: statement, proof, and the exact remaining gaps}
\author{}\date{}

\begin{document}\maketitle
\vspace{-1.2em}

\begin{abstract}\noindent
Note~7 (\texttt{universality\_normalization.tex}) states, as its Lemma~3.2, that the Bures and Kubo--Mori
second variations of a curve of normal states of a von Neumann algebra $\cM$ generated by a unitary flow are
given by the modular quadratic forms of the tangent vector $\eta=i(S^*-1)A\Om$, and leaves the lemma unproved
(``Proof route: Uhlmann's transition probability in standard form,
$\Fid(\omega,\omega_s)=\ip{\Om}{\Delta_{\omega_s,\omega}^{1/2}\Om}$, and Araki's perturbation theory\dots'').
That proof route is not available: the quoted formula for $\Fid$ is \emph{false}
(Section~\ref{sec:uhlmann}, Issue~\ref{iss:wrong-uhlmann}).  Here the \emph{fidelity} half of the lemma is proved
in full generality by a different route, and the Kubo--Mori half is proved in finite dimensions.
The key is a chain of three exact identities, valid for an arbitrary $\sigma$-finite $\cM$ in standard form
and an arbitrary vector $a=A\Om\in\Hil$ (Theorem~\ref{thm:three}):
\[
 \tfrac14\gB\;=\;\tfrac12\bigl\|(1-J)(1+\Delta)^{-1/2}a\bigr\|^2
 \;=\;\dist\bigl(a,\overline{\Mpsa\Om}\bigr)^2
 \;=\;\sup_{K\in\Msa}\bigl[\varphi(K)-\Var_\omega(K)\bigr].
\]
The middle expression is the ``distance to the gauge directions'', the right-hand one is a supremum over
$\cM$ that pairs with Alberti's variational formula for the transition probability.  The two ends of the chain
then give the two halves of the expansion:
$\Fid\ge1-\frac{s^2}8\gB+o(s^2)$ from the Uhlmann representation of $\Fid$ as a supremum over
$\U(\cM')$, and $\Fid\le1-\frac{s^2}8\gB+o(s^2)$ from Alberti's representation of $\Fid$ as an infimum over
positive invertible elements of $\cM$.  No hyperfiniteness, no type classification, no
approximation by finite-dimensional subalgebras, and no domain assumption beyond $\Om\in D(A)$ are needed.
The monotone approximation $\gB^{(n)}\uparrow\gB$ along $\cM_n\uparrow\cM$ is proved as a by-product
(Proposition~\ref{prop:mono}).  What remains open is stated precisely in Section~\ref{sec:open}:
(i) the Kubo--Mori half beyond finite dimensions; (ii) the passage from a finite-norm representative
$A\Om\in\Hil$ to the zero-collar modular-frame spectral integrals of Note~7, where $T(g)\Om\notin\Hil$;
(iii) norm-differentiability of the state curve at zero collar.
\end{abstract}

\tableofcontents

\section{What is proved here}\label{sec:overview}

\begin{center}\small
\begin{tabular}{@{}p{0.30\textwidth}p{0.62\textwidth}@{}}\toprule
\textbf{Item of the task} & \textbf{Status}\\\midrule
(A) tangent functional, Note~7 Lemma~3.1 & \textbf{proved} (Section~\ref{sec:tangent}); Note~7's proof omits the
domain condition $A\Om\in D(S^*)=D(\Delta^{-1/2})$ needed for $\eta$ to exist as a vector; the $J$-form of $\gB$ needs nothing.\\
(B) finite dimensions: $\gB$, $\gKM$, variational form & \textbf{proved} (Sections~\ref{sec:three}, \ref{sec:findim});
the constant $c$ may be dropped ($c\one\in\Mpsa$), and $c\in\mathbb C$ gives the same infimum \emph{because} $A=A^*$.\\
(C) Uhlmann representation of $\Fid$ & \textbf{proved} (Section~\ref{sec:uhlmann}), with the exact locator;
the formula quoted in Note~7 is refuted.\\
(D) lower bound $\Fid\ge1-\frac{s^2}8\gB+o(s^2)$ & \textbf{proved in general} (Section~\ref{sec:lower}), hypothesis: $\Om\in D(A)$ only.\\
(E) upper bound $\Fid\le1-\frac{s^2}8\gB+o(s^2)$ & \textbf{proved in general} (Section~\ref{sec:upper}), by Alberti's
infimum formula; the hyperfinite route of the task is \emph{not} needed and is shown to have a uniformity gap (Remark~\ref{rem:hyperfinite}).\\
(F) application to Note~7 & \textbf{proved for positive collar} $\eps>0$ (Section~\ref{sec:note7}), with $u_s=U(k_s)=e^{isT(g)}$
a strongly continuous one-parameter group; the zero-collar case is a genuine gap.\\
Kubo--Mori in general & \textbf{open} (Section~\ref{sec:open}, Gap~1).\\\bottomrule
\end{tabular}
\end{center}

\section{Setting, conventions, hypotheses}\label{sec:setting}

\textbf{Inner products} are conjugate-linear in the \emph{first} variable, as in Note~7.  For an antiunitary $J$
one has $\ip{J\xi}{J\zeta}=\ip{\zeta}{\xi}$ and hence
\begin{equation}\label{eq:Jswap}
 \ip{J\xi}{\zeta}=\ip{J\zeta}{\xi}\qquad(\xi,\zeta\in\Hil,\ J^2=1).
\end{equation}
\emph{Proof of \eqref{eq:Jswap}}: put $\alpha=\xi$, $\beta=J\zeta$ in $\ip{J\alpha}{J\beta}=\ip{\beta}{\alpha}$ to get
$\ip{J\xi}{\zeta}=\ip{J\xi}{J(J\zeta)}=\ip{J\zeta}{\xi}$.

\textbf{Standard form.}  $\cM$ is a von Neumann algebra with a cyclic and separating vector $\Om\in\Hil$;
$\omega=\ip{\Om}{\cdot\,\Om}$ is a faithful normal state.  $S_0:X\Om\mapsto X^*\Om$ ($X\in\cM$) and
$F_0:X'\Om\mapsto X'^*\Om$ ($X'\in\cM'$) are closable conjugate-linear operators; $S=\overline{S_0}$, $F=\overline{F_0}$,
$S^*=F$, $F^*=S$, $S=J\Delta^{1/2}$, $F=J\Delta^{-1/2}=\Delta^{1/2}J$, $J^2=1$, $J\Om=\Om$, $\Delta\Om=\Om$,
$J\cM J=\cM'$, $J\Delta J=\Delta^{-1}$ and therefore
\begin{equation}\label{eq:Jf}
 Jf(\Delta)J=\overline f(\Delta^{-1})\quad\text{for bounded Borel }f;\qquad\text{in particular}\qquad
 J(1+\Delta)^{-1/2}=\Delta^{1/2}(1+\Delta)^{-1/2}J .
\end{equation}
All of this is Tomita--Takesaki theory; the locator we use throughout is
O.~Bratteli, D.~W.~Robinson, \emph{Operator Algebras and Quantum Statistical Mechanics~1}, 2nd ed., Prop.~2.5.9
(closability, $S^*=F$, $F^*=S$) and Thm.~2.5.14 (Tomita's theorem, $J\cM J=\cM'$, $\Delta^{it}\cM\Delta^{-it}=\cM$).
The second identity in \eqref{eq:Jf} is the functional calculus of $J\Delta J=\Delta^{-1}$:
\begin{equation}\label{eq:Jf-proof}
 J(1+\Delta)^{-1/2}J=(1+\Delta^{-1})^{-1/2}=\bigl(\tfrac{\Delta}{1+\Delta}\bigr)^{1/2}=\Delta^{1/2}(1+\Delta)^{-1/2}.
\end{equation}
Throughout we abbreviate
\begin{equation}\label{eq:u-def}
 u:=(1+\Delta)^{-1/2},\qquad \norm{u}\le1,\qquad \norm{\Delta^{1/2}u}\le1,\qquad \norm{\Delta u^2}\le 1,
\end{equation}
all three bounded because $\lambda\mapsto(1+\lambda)^{-1/2},\lambda^{1/2}(1+\lambda)^{-1/2},\lambda(1+\lambda)^{-1}$
are bounded on $[0,\infty)$.  Equation~\eqref{eq:Jf} reads $Ju=\Delta^{1/2}uJ$, equivalently $JuJ=\Delta^{1/2}u$.

\textbf{Standard real subspaces.}
\begin{equation}\label{eq:realsub}
 \HM:=\overline{\Msa\Om},\qquad \HMp:=\overline{\Mpsa\Om}
\end{equation}
(closures in $\Hil$; these are closed \emph{real}-linear subspaces).  $\Msa$ and $\Mpsa$ denote the self-adjoint parts.

\textbf{The curve.}  We use two hypotheses, of decreasing strength.
\begin{itemize}[leftmargin=1.5em]
\item[(V)] \emph{(vector hypothesis)} There is a family $(\xi_s)_{|s|<s_0}$ of unit vectors in $\Hil$ with $\xi_0=\Om$,
$\omega_s:=\ip{\xi_s}{\cdot\,\xi_s}\big|_{\cM}$, which is differentiable in norm at $s=0$:
$\norm{\xi_s-\Om-s\dot\xi_0}=o(s)$.  We write $a:=i\dot\xi_0\in\Hil$ and call $a$ \emph{the tangent representative}.
\item[(F$\varphi$)] \emph{(functional hypothesis)} $(\omega_s)$ are normal states on $\cM$ and for every $y\in\cM$
$\omega_s(y)=\omega(y)+s\varphi(y)+o(s)$, with $\varphi(K)=-2\operatorname{Im}\ip{a}{K\Om}$ for $K\in\Msa$ and some $a\in\Hil$.
\end{itemize}
The model case is $\xi_s=u_s^*\Om=e^{-isA}\Om$ with $u_s=e^{isA}$ a strongly continuous unitary group on $\Hil$,
$A=A^*$ (in general affiliated with neither $\cM$ nor $\cM'$) and $\Om\in D(A)$: then $\dot\xi_0=-iA\Om$, so $a=A\Om$
(Lemma~\ref{lem:tangent}).  Note that (V) $\Rightarrow$ (F$\varphi$) with the same $a$ (Lemma~\ref{lem:tangent}\,(3)).
\emph{The vector $a$ is not unique}: only its class modulo $\HMp$ is (Corollary~\ref{cor:gauge}).

\textbf{Fidelity.}  $\Fid$ denotes \emph{root} fidelity, $\Fid(\rho,\sigma)=\Tr|\rho^{1/2}\sigma^{1/2}|$ in finite
dimensions, i.e.\ $\Fid=\sqrt{P}$ with $P$ the Uhlmann transition probability of Section~\ref{sec:uhlmann}.  This is the
convention of Notes~4,5,7.

\section{(A) The tangent functional}\label{sec:tangent}

\begin{lemma}[Differentiability and the tangent functional]\label{lem:tangent}
Let $A$ be self-adjoint on $\Hil$ with $\Om\in D(A)$, $u_s=e^{isA}$, $\xi_s=u_s^*\Om$, $\omega_s=\ip{\xi_s}{\cdot\,\xi_s}|_{\cM}$.
Then
\begin{enumerate}[leftmargin=1.7em]
\item $\norm{\xi_s-\Om+isA\Om}=o(s)$, i.e.\ hypothesis \textup{(V)} holds with $a=A\Om$;
\item for every $X\in\cM$,
 $\displaystyle\frac{d}{ds}\omega_s(X)\Big|_{s=0}=i\ip{A\Om}{X\Om}-i\ip{X^*\Om}{A\Om}=:\varphi(X)$,
 and if $A$ is bounded this equals $i\,\omega([A,X])$;
\item for $K\in\Msa$, $\varphi(K)=-2\operatorname{Im}\ip{A\Om}{K\Om}$; moreover $\varphi(\one)=0$ and
 $\operatorname{Im}\ip{\Om}{A\Om}=0$;
\item under hypothesis \textup{(V)} the same conclusions hold with $A\Om$ replaced by $a=i\dot\xi_0$.
\end{enumerate}
\end{lemma}

\lstep{1}{1}{$\lim_{s\to0}\norm{s^{-1}(e^{-isA}-1)\Om+iA\Om}=0$.}
\lby{Stone's theorem: for a self-adjoint $A$ and $\Om\in D(A)$, $s\mapsto e^{-isA}\Om$ is norm-differentiable with
derivative $-iAe^{-isA}\Om$ (Reed--Simon~I, Thm.~VIII.7(a): $\lim_{t\to0}t^{-1}(e^{-itA}-1)\Om=-iA\Om$ exactly for $\Om\in D(A)$)}
\lstep{1}{2}{Item (1) holds.}
\lby{$\langle1\rangle1$ and $\xi_s=e^{-isA}\Om$}
\lstep{1}{3}{$\omega_s(X)=\ip{\xi_s}{X\xi_s}$ is differentiable at $s=0$ with derivative $\ip{\dot\xi_0}{X\Om}+\ip{\Om}{X\dot\xi_0}$.}
\lproof
\lstep{2}{1}{$\omega_s(X)-\omega_0(X)=\ip{\xi_s-\Om}{X\xi_s}+\ip{\Om}{X(\xi_s-\Om)}$.}
\lby{bilinearity: $\ip{\xi_s}{X\xi_s}-\ip{\Om}{X\Om}=\ip{\xi_s-\Om}{X\xi_s}+\ip{\Om}{X\xi_s}-\ip{\Om}{X\Om}$}
\lstep{2}{2}{$s^{-1}\ip{\xi_s-\Om}{X\xi_s}\to\ip{\dot\xi_0}{X\Om}$ and $s^{-1}\ip{\Om}{X(\xi_s-\Om)}\to\ip{\Om}{X\dot\xi_0}$.}
\lby{$X$ is bounded, $s^{-1}(\xi_s-\Om)\to\dot\xi_0$ in norm by $\langle1\rangle1$, and $\xi_s\to\Om$ in norm;
continuity of the inner product in both arguments jointly on norm-convergent nets}
\lqed{2}{3}
\lby{$\langle2\rangle1$, $\langle2\rangle2$}
\lstep{1}{4}{With $\dot\xi_0=-iA\Om$: $\ip{\dot\xi_0}{X\Om}=i\ip{A\Om}{X\Om}$ and $\ip{\Om}{X\dot\xi_0}=-i\ip{X^*\Om}{A\Om}$.}
\lby{conjugate-linearity in the first slot gives $\ip{-iA\Om}{X\Om}=\overline{(-i)}\ip{A\Om}{X\Om}=i\ip{A\Om}{X\Om}$;
linearity in the second slot and $\ip{\Om}{XA\Om}=\ip{X^*\Om}{A\Om}$ give the second}
\lstep{1}{5}{Item (2) holds; and for bounded $A$, $\varphi(X)=i\ip{\Om}{A X\Om}-i\ip{\Om}{XA\Om}=i\,\omega([A,X])$.}
\lby{$\langle1\rangle3$, $\langle1\rangle4$; for bounded $A$, $\ip{A\Om}{X\Om}=\ip{\Om}{AX\Om}$ by self-adjointness}
\lstep{1}{6}{Item (3) holds.}
\lproof
\lstep{2}{1}{For $K=K^*$, $\ip{K^*\Om}{A\Om}=\ip{K\Om}{A\Om}=\overline{\ip{A\Om}{K\Om}}$.}
\lby{hermiticity of the inner product}
\lstep{2}{2}{$\varphi(K)=i\bigl(z-\bar z\bigr)$ with $z=\ip{A\Om}{K\Om}$, and $z-\bar z=2i\operatorname{Im}z$,
hence $\varphi(K)=-2\operatorname{Im}z$.}
\lby{$\langle1\rangle5$, $\langle2\rangle1$}
\lstep{2}{3}{$\ip{\Om}{A\Om}\in\mathbb R$, hence $\operatorname{Im}\ip{\Om}{A\Om}=0$ and $\varphi(\one)=0$.}
\lby{$A=A^*$ and $\Om\in D(A)$, so $\ip{\Om}{A\Om}=\ip{A\Om}{\Om}=\overline{\ip{\Om}{A\Om}}$; then $\langle2\rangle2$ with $K=\one$}
\lqed{2}{4}
\lby{$\langle2\rangle2$, $\langle2\rangle3$}
\lstep{1}{7}{Item (4) holds.}
\lproof
\lstep{2}{1}{$\langle1\rangle3$ used only $\norm{s^{-1}(\xi_s-\Om)-\dot\xi_0}\to0$, which is hypothesis (V).}
\lby{inspection of the proof of $\langle1\rangle3$}
\lstep{2}{2}{With $a:=i\dot\xi_0$, $\ip{\dot\xi_0}{X\Om}+\ip{\Om}{X\dot\xi_0}=i\ip{a}{X\Om}-i\ip{X^*\Om}{a}$,
and for $K=K^*$ this is $-2\operatorname{Im}\ip{a}{K\Om}$.}
\lby{$\dot\xi_0=-ia$ and the computation of $\langle1\rangle4$, $\langle1\rangle6$ verbatim}
\lstep{2}{3}{$\operatorname{Im}\ip{\Om}{a}=0$.}
\lby{$\norm{\xi_s}=1$ for all $s$ gives $0=\frac{d}{ds}\norm{\xi_s}^2|_0=2\operatorname{Re}\ip{\Om}{\dot\xi_0}
=2\operatorname{Re}\ip{\Om}{-ia}=-2\operatorname{Re}(i\ip{\Om}{a})=2\operatorname{Im}\ip{\Om}{a}$}
\lqed{2}{4}
\lby{$\langle2\rangle1$--$\langle2\rangle3$}
\lqed{1}{8}
\lby{$\langle1\rangle2$, $\langle1\rangle5$, $\langle1\rangle6$, $\langle1\rangle7$}

\begin{remark}\label{rem:selfadj-real}
Step $\langle2\rangle3$ of $\langle1\rangle6$ (equivalently $\langle2\rangle3$ of $\langle1\rangle7$) is the only place
where self-adjointness of $A$ (respectively normalisation $\norm{\xi_s}=1$) is used, and it is exactly the condition that
decides the question ``$c\in\mathbb R$ or $c\in\mathbb C$?'' of the task; see Proposition~\ref{prop:c-real}.
\end{remark}

\subsection{Note~7, Lemma~3.1 re-derived}\label{sec:eta}

\begin{lemma}[Note~7, Lemma~3.1, with its missing hypothesis]\label{lem:eta}
Assume the situation of Lemma~\ref{lem:tangent} and, in addition,
\begin{equation}\label{eq:Sdomain}
 A\Om\in D(S^*)=D(F)=D(\Delta^{-1/2}).
\end{equation}
Then $\varphi(X)=\ip{\eta}{X\Om}$ for all $X\in\cM$, with the unique vector
\begin{equation}\label{eq:eta}
 \eta=i(S^*-1)A\Om\in\Hil .
\end{equation}
Moreover:
\begin{enumerate}[leftmargin=1.7em]
\item $\eta$ is unchanged under $A\mapsto A+B'$ with $B'=B'^*\in\cM'$;
\item if $A\in\Msa$ and $A\Om\in D(\Delta)$ then $S^*A\Om=\Delta A\Om$ and $\eta=i(\Delta-1)A\Om$;
\item without any hypothesis beyond $A\Om\in\Hil$,
\begin{equation}\label{eq:gB-J}
 (1+\Delta)^{-1/2}(S^*-1)A\Om=(J-1)(1+\Delta)^{-1/2}A\Om ,
\end{equation}
both sides being everywhere defined and bounded in $A\Om$; in particular
$\gB:=2\norm{(1+\Delta)^{-1/2}\eta}^2=2\norm{(1-J)(1+\Delta)^{-1/2}A\Om}^2\le8\norm{A\Om}^2$.
\end{enumerate}
Conversely, if $A\Om\notin D(S^*)$ then there is \emph{no} vector $\eta$ with $\varphi=\ip{\eta}{\cdot\,\Om}$.
\end{lemma}

\lstep{1}{1}{For $X\in\cM$: $\ip{X^*\Om}{A\Om}=\ip{SX\Om}{A\Om}$.}
\lby{$S X\Om=X^*\Om$ for $X\in\cM$, by definition of $S_0$ and $S\supseteq S_0$}
\lstep{1}{2}{$\ip{SX\Om}{A\Om}=\ip{S^*A\Om}{X\Om}$, and this is where \eqref{eq:Sdomain} is used.}
\lby{definition of the adjoint of a densely defined conjugate-linear operator: $\ip{S^*\zeta}{\xi}=\ip{S\xi}{\zeta}$ for
$\xi\in D(S),\zeta\in D(S^*)$; here $\xi=X\Om$, $\zeta=A\Om$, legitimate by \eqref{eq:Sdomain}}
\lstep{1}{3}{$\varphi(X)=i\ip{A\Om}{X\Om}-i\ip{S^*A\Om}{X\Om}=\ip{i(S^*-1)A\Om}{X\Om}$.}
\lby{Lemma~\ref{lem:tangent}(2), $\langle1\rangle1$, $\langle1\rangle2$, and $i\ip{\zeta}{\cdot}=\ip{-i\zeta}{\cdot}$,
$-i\ip{\zeta}{\cdot}=\ip{i\zeta}{\cdot}$, so $i\ip{A\Om}{\cdot}-i\ip{S^*A\Om}{\cdot}=\ip{-iA\Om+iS^*A\Om}{\cdot}$}
\lstep{1}{4}{$\eta$ is unique.}
\lby{$\cM\Om$ is dense ($\Om$ cyclic), and a vector is determined by its inner products with a dense set}
\lstep{1}{5}{Item (1).}
\lby{$S^*B'\Om=FB'\Om=B'^*\Om=B'\Om$ for $B'\in\Mpsa$ (definition of $F_0$ and $F=\overline{F_0}\supseteq F_0$), so
$(S^*-1)B'\Om=0$; and $D(S^*)$ is a linear space containing $\cM'\Om$}
\lstep{1}{6}{Item (2).}
\lproof
\lstep{2}{1}{For $X\in\cM$: $JX\Om=\Delta^{1/2}X^*\Om$.}
\lby{apply $J$ to $J\Delta^{1/2}X^*\Om=SX^*\Om=X\Om$ and use $J^2=1$}
\lstep{2}{2}{For $A\in\Msa$: $JA\Om=\Delta^{1/2}A\Om$, and $S^*A\Om=\Delta^{1/2}JA\Om=\Delta^{1/2}\Delta^{1/2}A\Om=\Delta A\Om$,
the last step requiring $\Delta^{1/2}A\Om\in D(\Delta^{1/2})$, i.e.\ $A\Om\in D(\Delta)$.}
\lby{$\langle2\rangle1$ with $X=A=A^*$; $S^*=F=\Delta^{1/2}J$; and $\Delta^{1/2}\Delta^{1/2}\subseteq\Delta$ with
$D(\Delta)=\{\zeta\in D(\Delta^{1/2}):\Delta^{1/2}\zeta\in D(\Delta^{1/2})\}$ by the spectral theorem}
\lqed{2}{3}
\lby{$\langle2\rangle2$ and \eqref{eq:eta}}
\lstep{1}{7}{Item (3).}
\lproof
\lstep{2}{1}{$u\,S^*=u\Delta^{1/2}J=Ju$ as everywhere-defined bounded conjugate-linear operators, $u=(1+\Delta)^{-1/2}$.}
\lby{\eqref{eq:Jf}: $Ju=\Delta^{1/2}uJ=u\Delta^{1/2}J$, the last equality because $u$ and $\Delta^{1/2}$ commute
(both are functions of $\Delta$) and $u\Delta^{1/2}=\Delta^{1/2}u$ is bounded by \eqref{eq:u-def}}
\lstep{2}{2}{If $A\Om\in D(S^*)$ then $u(S^*-1)A\Om=(J-1)uA\Om$.}
\lby{$\langle2\rangle1$ applied to $A\Om$, and $u$ linear}
\lstep{2}{3}{The right-hand side of \eqref{eq:gB-J} is defined for every $A\Om\in\Hil$ and
$\norm{(J-1)uA\Om}\le2\norm{A\Om}$.}
\lby{$J$ antiunitary and $\norm u\le1$}
\lqed{2}{4}
\lby{$\langle2\rangle2$, $\langle2\rangle3$; $\gB=2\norm{u\eta}^2=2\norm{u\,i(S^*-1)A\Om}^2=2\norm{(1-J)uA\Om}^2$
because $|i|=1$ and $\norm{(J-1)\zeta}=\norm{(1-J)\zeta}$}
\lstep{1}{8}{The converse.}
\lby{if $\varphi=\ip{\eta}{\cdot\,\Om}$ then by $\langle1\rangle3$'s first equality
$\ip{SX\Om}{A\Om}=i^{-1}\bigl(i\ip{A\Om}{X\Om}-\varphi(X)\bigr)$ is bounded by $\mathrm{const}\cdot\norm{X\Om}$
uniformly in $X\in\cM$, which is precisely the statement $A\Om\in D(S^*)$ (definition of the adjoint of the
densely defined operator $S$)}
\lqed{1}{9}
\lby{$\langle1\rangle3$--$\langle1\rangle8$}

\begin{verdictgap}
Note~7, Lemma~3.1 and its proof (``$\ip{X^*\Om}{A\Om}=\ip{J\Delta^{1/2}X\Om}{A\Om}=\ip{\Delta^{1/2}JA\Om}{X\Om}
=\ip{S^*A\Om}{X\Om}$, which gives $\eta$'') is correct as an identity but omits the hypothesis
$A\Om\in D(S^*)=D(\Delta^{-1/2})$, which is automatic in finite dimensions and is a genuine restriction in general:
the middle equality moves $\Delta^{1/2}$ across the inner product and needs $JA\Om\in D(\Delta^{1/2})$.
Likewise, item (2) of Note~7's Lemma~3.1 (``If $A\in\cM$, then $S^*A\Om=\Delta A\Om$'') needs $A\Om\in D(\Delta)$,
which is \emph{not} automatic for bounded $A\in\cM$ ($\cM\Om\subseteq D(\Delta^{1/2})$ always, $\cM\Om\subseteq D(\Delta)$
almost never).  Item~(3) of Note~7's Lemma~3.1, the formula
$\gB=2\norm{(1-J)(1+\Delta)^{-1/2}A\Om}^2$, is correct with no hypothesis at all, and is the form that should be used
as the definition; this is Lemma~\ref{lem:eta}(3).  No numerical or structural conclusion of Note~7 is affected.
\end{verdictgap}

\section{The three faces of the Bures form}\label{sec:three}

This section is the technical core.  Everything in it is exact: no expansion, no limit, no hypothesis on $\cM$
beyond having a cyclic separating vector, and no hypothesis on $a\in\Hil$ at all.

\begin{definition}\label{def:gB}
For $a\in\Hil$ put $\displaystyle \gB(a):=2\bigl\|(1-J)(1+\Delta)^{-1/2}a\bigr\|^2$ and
$\Var_\omega(K):=\omega(K^2)-\omega(K)^2$ for $K\in\Msa$.  By Lemma~\ref{lem:eta}(3) this agrees with
$2\norm{(1+\Delta)^{-1/2}\eta}^2=2\ip{\eta}{(1+\Delta)^{-1}\eta}$ whenever $\eta=i(S^*-1)a$ exists.
\end{definition}

\begin{lemma}[The two standard real subspaces]\label{lem:realsub}
$\HM=\{\xi\in D(S):S\xi=\xi\}$ and $\HMp=\{\xi\in D(F):F\xi=\xi\}$.
\end{lemma}

\lstep{1}{1}{$\Msa\Om\subseteq\{\xi\in D(S):S\xi=\xi\}$, and the latter set is closed.}
\lby{$SK\Om=K^*\Om=K\Om$ for $K\in\Msa$; the set is closed because $S$ is a closed operator: if $\xi_n\to\xi$ with
$S\xi_n=\xi_n\to\xi$ then $(\xi,\xi)$ is in the graph of $S$}
\lstep{1}{2}{$\HM\subseteq\{\xi\in D(S):S\xi=\xi\}$.}
\lby{$\langle1\rangle1$ and the definition of $\HM$ as a closure}
\lstep{1}{3}{Conversely, let $\xi\in D(S)$ with $S\xi=\xi$.  Then $\xi\in\HM$.}
\lproof
\lstep{2}{1}{$D(S)$ is the closure of $\cM\Om$ in the graph norm, so there are $X_n\in\cM$ with
$X_n\Om\to\xi$ and $X_n^*\Om=SX_n\Om\to S\xi=\xi$.}
\lby{$S=\overline{S_0}$ with $D(S_0)=\cM\Om$}
\lstep{2}{2}{$K_n:=\tfrac12(X_n+X_n^*)\in\Msa$ and $K_n\Om=\tfrac12(X_n\Om+X_n^*\Om)\to\tfrac12(\xi+\xi)=\xi$.}
\lby{$\langle2\rangle1$}
\lqed{2}{3}
\lby{$\langle2\rangle2$ and $\HM=\overline{\Msa\Om}$}
\lstep{1}{4}{The statement for $\cM'$.}
\lby{$F=\overline{F_0}$ with $D(F_0)=\cM'\Om$ (Bratteli--Robinson Prop.~2.5.9), so $\langle1\rangle1$--$\langle1\rangle3$
apply verbatim with $(\cM,S)$ replaced by $(\cM',F)$}
\lqed{1}{5}
\lby{$\langle1\rangle2$, $\langle1\rangle3$, $\langle1\rangle4$}

\begin{lemma}[Explicit real-orthogonal projection onto $\HMp$]\label{lem:proj}
Write
\begin{equation}\label{eq:uvm}
 u=(1+\Delta)^{-1/2},\qquad v=\Delta^{1/2}(1+\Delta)^{-1/2},\qquad m=uv=\Delta^{1/2}(1+\Delta)^{-1},
\end{equation}
bounded, self-adjoint, positive, with
\begin{equation}\label{eq:uvm-rel}
 u^2+v^2=1,\qquad Ju=vJ,\qquad Jv=uJ,\qquad Jm=mJ,\qquad \Delta^{1/2}u^2=m,\qquad \Delta^{1/2}m=v^2 .
\end{equation}
For $b\in\Hil$ set
\begin{equation}\label{eq:Eprime}
 E'b:=v^2b+mJb\;=\;\Delta(1+\Delta)^{-1}b+(1+\Delta)^{-1}\Delta^{1/2}Jb .
\end{equation}
Then $E'b\in\HMp$; $b-E'b=u(1-J)ub$; $\operatorname{Re}\ip{b-E'b}{\zeta}=0$ for every $\zeta\in\HMp$; and
\begin{equation}\label{eq:dist-formula}
 \dist(b,\HMp)^2=\norm{b-E'b}^2=\tfrac12\bigl\|(1-J)ub\bigr\|^2
 =\norm{ub}^2-\operatorname{Re}\ip{Jb}{mb}=\tfrac14\gB(b).
\end{equation}
Applying the same statement to the standard form $(\cM',\Hil,J,\Delta^{-1},\Om)$ of the commutant,
\begin{equation}\label{eq:dist-HM}
 \dist(b,\HM)^2=\tfrac12\norm{(1-J)vb}^2=\norm{vb}^2-\operatorname{Re}\ip{Jb}{mb}.
\end{equation}
\end{lemma}

\lstep{1}{1}{The relations \eqref{eq:uvm-rel} hold and $u,v,m$ are bounded.}
\lby{$u=f_1(\Delta),v=f_2(\Delta),m=f_3(\Delta)$ with the bounded continuous real functions
$f_1(\lambda)=(1+\lambda)^{-1/2}$, $f_2(\lambda)=\lambda^{1/2}(1+\lambda)^{-1/2}$,
$f_3(\lambda)=\lambda^{1/2}(1+\lambda)^{-1}\le\tfrac12$;
$f_1^2+f_2^2=1$; $Jf(\Delta)J=f(\Delta^{-1})$ by \eqref{eq:Jf}, and
$f_1(\lambda^{-1})=f_2(\lambda)$, $f_2(\lambda^{-1})=f_1(\lambda)$, $f_3(\lambda^{-1})=f_3(\lambda)$;
$\lambda^{1/2}f_1(\lambda)^2=f_3(\lambda)$ and $\lambda^{1/2}f_3(\lambda)=f_2(\lambda)^2$}
\lstep{1}{2}{$b-E'b=u(1-J)ub$.}
\lby{$b-v^2b=u^2b$ by \eqref{eq:uvm-rel}; $mJb=uvJb=u(Ju)b$ by $vJ=Ju$; hence
$b-E'b=u^2b-uJub=u(ub-Jub)=u(1-J)ub$}
\lstep{1}{3}{$JE'b=u^2Jb+mb$.}
\lby{$Jv^2=(Jv)v=(uJ)v=u(Jv)=u(uJ)=u^2J$, so $Jv^2b=u^2Jb$; and $JmJb=mJJb=mb$ by $Jm=mJ$, $J^2=1$}
\lstep{1}{4}{$JE'b\in D(\Delta^{1/2})$ and $\Delta^{1/2}JE'b=mJb+v^2b=E'b$; hence $E'b\in D(F)$ and $FE'b=E'b$.}
\lproof
\lstep{2}{1}{$u^2Jb\in D(\Delta^{1/2})$ and $\Delta^{1/2}u^2Jb=mJb$.}
\lby{$\int_0^\infty\lambda\,f_1(\lambda)^4\,d\norm{E_\lambda Jb}^2\le\int_0^\infty\lambda(1+\lambda)^{-2}
 d\norm{E_\lambda Jb}^2\le\tfrac14\norm{b}^2<\infty$ ($E_\lambda$ the spectral resolution of $\Delta$), which is the
 spectral criterion for $u^2Jb\in D(\Delta^{1/2})$; the value is $\Delta^{1/2}u^2Jb=f_3(\Delta)Jb=mJb$
 by \eqref{eq:uvm-rel}}
\lstep{2}{2}{$mb\in D(\Delta^{1/2})$ and $\Delta^{1/2}mb=v^2b$.}
\lby{$\int_0^\infty\lambda f_3(\lambda)^2d\norm{E_\lambda b}^2=\int_0^\infty\lambda^2(1+\lambda)^{-2}
 d\norm{E_\lambda b}^2\le\norm b^2<\infty$; the value is $f_2(\Delta)^2b=v^2b$ by \eqref{eq:uvm-rel}}
\lstep{2}{3}{$\Delta^{1/2}JE'b=mJb+v^2b=E'b$.}
\lby{$\langle1\rangle3$, $\langle2\rangle1$, $\langle2\rangle2$, additivity of $\Delta^{1/2}$ on its domain, and
$E'b=v^2b+mJb$ from \eqref{eq:Eprime}}
\lqed{2}{4}
\lby{$F=\Delta^{1/2}J$ and $\langle2\rangle3$}
\lstep{1}{5}{$E'b\in\HMp$.}
\lby{$\langle1\rangle4$ and Lemma~\ref{lem:realsub}}
\lstep{1}{6}{For $\zeta\in\HMp$: $Ju\zeta=u\zeta$.}
\lproof
\lstep{2}{1}{$\zeta\in D(\Delta^{-1/2})$ and $J\zeta=\Delta^{-1/2}\zeta$.}
\lby{Lemma~\ref{lem:realsub}: $\Delta^{1/2}J\zeta=F\zeta=\zeta$, so $J\zeta\in D(\Delta^{1/2})$ and, applying
$\Delta^{-1/2}$ (defined on $\operatorname{Ran}\Delta^{1/2}$ up to the kernel of $\Delta$, which is $0$),
$J\zeta=\Delta^{-1/2}\zeta$; in particular $\zeta\in D(\Delta^{-1/2})$}
\lstep{2}{2}{$Ju\zeta=vJ\zeta=v\Delta^{-1/2}\zeta=u\zeta$.}
\lby{$Ju=vJ$; $\langle2\rangle1$; and
$v\Delta^{-1/2}\zeta=f_2(\Delta)\Delta^{-1/2}\zeta=f_1(\Delta)\zeta=u\zeta$ because
$f_2(\lambda)\lambda^{-1/2}=f_1(\lambda)$ and $\zeta\in D(\Delta^{-1/2})$}
\lqed{2}{3}
\lby{$\langle2\rangle2$}
\lstep{1}{7}{$\ip{b-E'b}{\zeta}$ is purely imaginary for $\zeta\in\HMp$; hence $\operatorname{Re}\ip{b-E'b}{\zeta}=0$.}
\lproof
\lstep{2}{1}{$\ip{b-E'b}{\zeta}=\ip{ub}{u\zeta}-\ip{Jub}{u\zeta}$.}
\lby{$\langle1\rangle2$ and $u=u^*$}
\lstep{2}{2}{$\ip{Jub}{u\zeta}=\ip{Ju\zeta}{ub}=\ip{u\zeta}{ub}=\overline{\ip{ub}{u\zeta}}$.}
\lby{\eqref{eq:Jswap}; $\langle1\rangle6$; hermiticity of the inner product}
\lqed{2}{3}
\lby{$\langle2\rangle1$, $\langle2\rangle2$: $z-\bar z\in i\mathbb R$}
\lstep{1}{8}{$\dist(b,\HMp)^2=\norm{b-E'b}^2$.}
\lby{for $\zeta\in\HMp$, $\norm{b-\zeta}^2=\norm{b-E'b}^2+\norm{E'b-\zeta}^2+2\operatorname{Re}\ip{b-E'b}{E'b-\zeta}$,
and the last term vanishes by $\langle1\rangle5$, $\langle1\rangle7$ since $E'b-\zeta\in\HMp$ ($\HMp$ is a real subspace);
so $\norm{b-\zeta}\ge\norm{b-E'b}$ with equality at $\zeta=E'b$}
\lstep{1}{9}{$\norm{u(1-J)ub}^2=\tfrac12\norm{(1-J)ub}^2$.}
\lproof
\lstep{2}{1}{$\psi:=(1-J)ub$ satisfies $J\psi=-\psi$.}
\lby{$J\psi=Jub-J^2ub=Jub-ub=-\psi$}
\lstep{2}{2}{$\ip{\psi}{u^2\psi}=\ip{\psi}{v^2\psi}$.}
\lby{$\ip{u^2\psi}{\psi}=\ip{J\psi}{Ju^2\psi}$ by $\ip{J\alpha}{J\beta}=\ip{\beta}{\alpha}$ with
$\alpha=\psi,\beta=u^2\psi$; $Ju^2=v^2J$ by $\langle1\rangle3$'s computation; so
$\ip{u^2\psi}{\psi}=\ip{-\psi}{v^2J\psi}=\ip{-\psi}{v^2(-\psi)}=\ip{\psi}{v^2\psi}$; and
$\ip{u^2\psi}{\psi}=\ip{\psi}{u^2\psi}$ because $u^2\ge0$}
\lstep{2}{3}{$2\ip{\psi}{u^2\psi}=\ip{\psi}{(u^2+v^2)\psi}=\norm\psi^2$.}
\lby{$\langle2\rangle2$ and $u^2+v^2=1$}
\lqed{2}{4}
\lby{$\norm{u\psi}^2=\ip{\psi}{u^2\psi}$ and $\langle2\rangle3$}
\lstep{1}{10}{$\tfrac12\norm{(1-J)ub}^2=\norm{ub}^2-\operatorname{Re}\ip{Jb}{mb}$.}
\lby{$\norm{(1-J)ub}^2=\norm{ub}^2-2\operatorname{Re}\ip{Jub}{ub}+\norm{Jub}^2=2\norm{ub}^2-2\operatorname{Re}\ip{Jub}{ub}$
($J$ antiunitary), and $\ip{Jub}{ub}=\ip{vJb}{ub}=\ip{Jb}{vub}=\ip{Jb}{mb}$ by $Ju=vJ$ and $v=v^*$}
\lstep{1}{11}{Equation \eqref{eq:dist-HM}.}
\lby{$(\cM',\Hil,J,\Om)$ is a standard form with modular operator $\Delta'=\Delta^{-1}$, modular conjugation $J'=J$ and
Tomita operator $S'=F$, $F'=S$ (Bratteli--Robinson Prop.~2.5.9 and Thm.~2.5.14); therefore
$u'=(1+\Delta')^{-1/2}=f_1(\Delta^{-1})=f_2(\Delta)=v$, $m'=f_3(\Delta^{-1})=m$, and
$\mathsf H_{(\cM')'}=\HM$; substituting into \eqref{eq:dist-formula}}
\lqed{1}{12}
\lby{$\langle1\rangle5$, $\langle1\rangle7$--$\langle1\rangle11$ and Definition~\ref{def:gB}}

\begin{theorem}[Three faces of the Bures form]\label{thm:three}
Let $\cM$ be a von Neumann algebra with cyclic separating $\Om$, $a\in\Hil$ with $\operatorname{Im}\ip{\Om}{a}=0$,
and $\varphi(K):=-2\operatorname{Im}\ip{a}{K\Om}$ for $K\in\Msa$.  Then
\begin{equation}\label{eq:three}
 \tfrac14\gB(a)\;=\;\tfrac12\bigl\|(1-J)(1+\Delta)^{-1/2}a\bigr\|^2
 \;=\;\dist\bigl(a,\overline{\Mpsa\Om}\bigr)^2
 \;=\;\sup_{K\in\Msa}\bigl[\varphi(K)-\Var_\omega(K)\bigr].
\end{equation}
The infimum in the third expression is attained at $E'a$ of \eqref{eq:Eprime}; the supremum in the fourth is attained
in the limit along any sequence $K_n\in\Msa$ with $K_n\Om\to E_{\cM}(-ia)$, $E_{\cM}$ the real-orthogonal projection
onto $\HM$.  In particular the right-hand side depends on $a$ only through $\varphi$.
\end{theorem}

\lstep{1}{1}{The first two equalities of \eqref{eq:three}.}
\lby{Definition~\ref{def:gB} and Lemma~\ref{lem:proj}, equation \eqref{eq:dist-formula}, with $b=a$}
\lstep{1}{2}{\ldefine $w:=-ia$.  Then $\varphi(K)=2\operatorname{Re}\ip{w}{K\Om}$ for $K\in\Msa$, and
$\operatorname{Re}\ip{w}{\Om}=0$.}
\lby{$\ip{-ia}{k}=\overline{(-i)}\ip{a}{k}=i\ip{a}{k}$, so $\operatorname{Re}\ip{w}{k}=\operatorname{Re}(i\ip{a}{k})
=-\operatorname{Im}\ip{a}{k}$, hence $2\operatorname{Re}\ip{w}{K\Om}=-2\operatorname{Im}\ip{a}{K\Om}=\varphi(K)$;
and $\operatorname{Re}\ip{w}{\Om}=-\operatorname{Im}\ip{a}{\Om}=\operatorname{Im}\ip{\Om}{a}=0$ by hypothesis}
\lstep{1}{3}{\lsuffices \lprove $\displaystyle\sup_{K\in\Msa}\bigl[\varphi(K)-\Var_\omega(K)\bigr]
 =\sup_{k\in\HM}\bigl[2\operatorname{Re}\ip{w}{k}-\norm k^2\bigr]$ and
 $\displaystyle\sup_{k\in\HM}\bigl[2\operatorname{Re}\ip{w}{k}-\norm k^2\bigr]=\dist(a,\HMp)^2$.}
\lby{$\langle1\rangle1$ and transitivity of equality}
\lstep{1}{4}{For $K\in\Msa$ and $c=\omega(K)\in\mathbb R$: $\varphi(K-c\one)=\varphi(K)$ and
$\Var_\omega(K-c\one)=\Var_\omega(K)$; and for $K$ with $\omega(K)=0$,
$\varphi(K)-\Var_\omega(K)=2\operatorname{Re}\ip{w}{K\Om}-\norm{K\Om}^2$.}
\lby{$\varphi(\one)=-2\operatorname{Im}\ip{a}{\Om}=0$ by hypothesis and $\varphi$ real-linear on $\Msa$;
$\Var_\omega$ is invariant under adding real multiples of $\one$ by its definition;
$\Var_\omega(K)=\omega(K^2)-\omega(K)^2=\norm{K\Om}^2$ when $\omega(K)=0$}
\lstep{1}{5}{$\displaystyle\sup_{K\in\Msa}[\varphi(K)-\Var_\omega(K)]
 =\sup\bigl\{2\operatorname{Re}\ip{w}{k}-\norm k^2:\ k\in L\bigr\}$, where
 $L:=\{k\in\HM:\operatorname{Re}\ip{\Om}{k}=0\}$.}
\lproof
\lstep{2}{1}{$\le$: every $K\in\Msa$ may be replaced by $K-\omega(K)\one\in\Msa$ without changing
$\varphi(K)-\Var_\omega(K)$, and then $k=K\Om\in L$.}
\lby{$\langle1\rangle4$; $K\Om\in\Msa\Om\subseteq\HM$ and $\operatorname{Re}\ip{\Om}{K\Om}=\omega(K)=0$}
\lstep{2}{2}{$\ge$: $\{K\Om:K\in\Msa,\ \omega(K)=0\}$ is dense in $L$.}
\lby{given $k\in L$, pick $K_n\in\Msa$ with $K_n\Om\to k$ ($\HM=\overline{\Msa\Om}$); then
$\omega(K_n)=\ip{\Om}{K_n\Om}\to\ip{\Om}{k}$, which is real (a limit of reals) and has vanishing real part, hence is $0$;
so $\widetilde K_n:=K_n-\omega(K_n)\one\in\Msa$ satisfies $\omega(\widetilde K_n)=0$ and
$\widetilde K_n\Om=K_n\Om-\omega(K_n)\Om\to k$}
\lstep{2}{3}{$k\mapsto2\operatorname{Re}\ip{w}{k}-\norm k^2$ is continuous on $\Hil$.}
\lby{continuity of the inner product and of the norm}
\lqed{2}{4}
\lby{$\langle2\rangle1$--$\langle2\rangle3$: a supremum over a dense subset of a continuous function equals the
supremum over the closure}
\lstep{1}{6}{For any closed real-linear subspace $R\subseteq\Hil$ and any $w\in\Hil$,
$\sup_{k\in R}[2\operatorname{Re}\ip{w}{k}-\norm k^2]=\norm{P_Rw}^2=\norm w^2-\dist(w,R)^2$,
with $P_R$ the orthogonal projection of the real Hilbert space $(\Hil,\operatorname{Re}\ip\cdot\cdot)$ onto $R$.}
\lby{$(\Hil,\operatorname{Re}\ip\cdot\cdot)$ is a real Hilbert space and $R$ a closed subspace, so $P_R$ exists and
$\operatorname{Re}\ip{w}{k}=\operatorname{Re}\ip{P_Rw}{k}$ for $k\in R$; then
$2\operatorname{Re}\ip{w}{k}-\norm k^2=\norm{P_Rw}^2-\norm{k-P_Rw}^2$, maximal exactly at $k=P_Rw$;
Pythagoras gives $\norm{P_Rw}^2=\norm w^2-\norm{w-P_Rw}^2$}
\lstep{1}{7}{$\sup_{k\in L}[\cdots]=\sup_{k\in\HM}[\cdots]=\norm{w}^2-\dist(w,\HM)^2$.}
\lby{$\HM=L\oplus_{\mathbb R}\mathbb R\Om$ orthogonally in $\operatorname{Re}\ip\cdot\cdot$
($\Om=\one\cdot\Om\in\Msa\Om\subseteq\HM$, $\norm\Om=1$, and $L$ is by definition the part of $\HM$ real-orthogonal
to $\Om$), so $P_{\HM}w=P_Lw+\operatorname{Re}\ip{\Om}{w}\,\Om=P_Lw$ by $\langle1\rangle2$; then $\langle1\rangle6$
twice}
\lstep{1}{8}{$\norm w^2-\dist(w,\HM)^2=\norm{ua}^2-\operatorname{Re}\ip{Ja}{ma}$.}
\lproof
\lstep{2}{1}{$\dist(w,\HM)^2=\norm{vw}^2-\operatorname{Re}\ip{Jw}{mw}$.}
\lby{\eqref{eq:dist-HM} with $b=w$}
\lstep{2}{2}{$\norm w^2-\norm{vw}^2=\norm{uw}^2=\norm{ua}^2$.}
\lby{$u^2+v^2=1$ gives $\norm{uw}^2+\norm{vw}^2=\ip{w}{(u^2+v^2)w}=\norm w^2$; and $\norm{uw}=\norm{u(-ia)}=\norm{ua}$}
\lstep{2}{3}{$\operatorname{Re}\ip{Jw}{mw}=-\operatorname{Re}\ip{Ja}{ma}$.}
\lby{$Jw=J(-ia)=\overline{(-i)}Ja=iJa$ ($J$ conjugate-linear) and $mw=-ima$, so
$\ip{Jw}{mw}=\ip{iJa}{-ima}=\overline{i}\,(-i)\ip{Ja}{ma}=(-i)(-i)\ip{Ja}{ma}=-\ip{Ja}{ma}$}
\lqed{2}{4}
\lby{$\langle2\rangle1$--$\langle2\rangle3$}
\lstep{1}{9}{$\norm{ua}^2-\operatorname{Re}\ip{Ja}{ma}=\dist(a,\HMp)^2$.}
\lby{\eqref{eq:dist-formula} with $b=a$}
\lqed{1}{10}
\lby{$\langle1\rangle3$, $\langle1\rangle5$, $\langle1\rangle7$, $\langle1\rangle8$, $\langle1\rangle9$}

\begin{verdictok}
All four expressions in \eqref{eq:three} were checked numerically against each other, and against the SLD Fisher
information $2\sum_{ij}|d\rho_{ij}|^2/(p_i+p_j)$ and against $2\ip{\eta}{(1+\Delta)^{-1}\eta}$, for a random faithful
$\rho$ on $\mathbb C^4$ and a random self-adjoint $A$ on the $16$-dimensional standard-form space: all six numbers agree
to $10$ significant digits (script \texttt{scratchpad/sv\_check.py}; values
$20.442477994614876,\ 20.442477994614897,\ 20.442477994698837,\ 20.44247799459117,\ 20.442477994614915$).
\end{verdictok}

\begin{corollary}[Gauge invariance; the role of $c$]\label{cor:gauge}
\begin{enumerate}[leftmargin=1.7em]
\item $\gB(a)$ depends only on the functional $\varphi$, not on the representative $a$; in particular
$\gB(a)=\gB(a+\zeta)$ for every $\zeta\in\HMp$, so $A$ may be changed by any $B'\in\Mpsa$ (and by any real multiple of
$\one$, since $\one\in\Mpsa$).
\item $\dist(a,\HMp)^2=\inf\{\norm{a-B'\Om-c\Om}^2:B'\in\Mpsa,\ c\in\mathbb R\}$: the constant $c$ is redundant.
\end{enumerate}
\end{corollary}
\lstep{1}{1}{(1).}
\lby{by Theorem~\ref{thm:three} $\gB(a)=4\sup_{K\in\Msa}[\varphi(K)-\Var_\omega(K)]$, an expression in $\varphi$ alone;
and if $\zeta\in\HMp$ then $\varphi$ is unchanged, because for $\zeta=B'\Om$ with $B'\in\Mpsa$ and $K\in\Msa$,
$\operatorname{Im}\ip{B'\Om}{K\Om}=\operatorname{Im}\ip{\Om}{B'K\Om}=\operatorname{Im}\ip{\Om}{KB'\Om}
=\operatorname{Im}\ip{K\Om}{B'\Om}=-\operatorname{Im}\ip{B'\Om}{K\Om}$, hence $=0$, and this extends to
$\zeta\in\HMp$ by continuity; also $\operatorname{Im}\ip{\Om}{a+\zeta}=\operatorname{Im}\ip{\Om}{a}=0$ since
$\ip{\Om}{B'\Om}\in\mathbb R$}
\lstep{1}{2}{(2).}
\lby{$\one\in\Mpsa$, so $c\Om=(c\one)\Om\in\Mpsa\Om$ and $B'\Om+c\Om=(B'+c\one)\Om\in\Mpsa\Om$; conversely every
$B'\Om$ arises with $c=0$; hence the two infima are over the same set $\Mpsa\Om$, whose closure is $\HMp$, and
$\inf$ over a set equals $\inf$ over its closure for the continuous function $\norm{a-\cdot}^2$}

\begin{proposition}[$c$ real or complex?]\label{prop:c-real}
Let $G_{\mathbb R}:=\inf\{\norm{a-B'\Om-c\Om}^2:B'\in\Mpsa,c\in\mathbb R\}$ and $G_{\mathbb C}$ the same infimum with
$c\in\mathbb C$.  Then $G_{\mathbb C}=G_{\mathbb R}-\bigl(\operatorname{Im}\ip{\Om}{a}\bigr)^2$.
Consequently $G_{\mathbb C}=G_{\mathbb R}=\tfrac14\gB(a)$ \emph{if and only if} $\operatorname{Im}\ip{\Om}{a}=0$,
which holds automatically when $a=A\Om$ with $A=A^*$ and $\Om\in D(A)$
\textup{(}Lemma~\ref{lem:tangent}\,(3)\textup{)} and, more generally, under hypothesis \textup{(V)}
\textup{(}Lemma~\ref{lem:tangent}\,(4)\textup{)}.  The correct range is therefore $c\in\mathbb R$;
allowing $c\in\mathbb C$ is harmless \emph{only because} $A$ is self-adjoint, and would strictly lower the infimum
for a general $a$ \textup{(}e.g.\ $a=i\Om$: $G_{\mathbb R}=1$, $G_{\mathbb C}=0$, while $\gB(i\Om)=4$\textup{)}.
\end{proposition}
\lstep{1}{1}{$G_{\mathbb C}=\inf\{\norm{a-\zeta-it\Om}^2:\zeta\in\HMp,\ t\in\mathbb R\}$.}
\lby{writing $c=c_1+ic_2$ and absorbing $c_1\one\in\Mpsa$ into $B'$ (Corollary~\ref{cor:gauge}(2)), and passing to
closures as there}
\lstep{1}{2}{$\operatorname{Re}\ip{a-E'a}{i\Om}=\operatorname{Im}\ip{\Om}{a}$ and $i\Om\perp_{\mathbb R}\HMp$.}
\lproof
\lstep{2}{1}{$\operatorname{Re}\ip{i\Om}{\zeta}=\operatorname{Im}\ip{\Om}{\zeta}=0$ for $\zeta\in\Mpsa\Om$, hence for
$\zeta\in\HMp$.}
\lby{$\ip{\Om}{B'\Om}=\omega'(B')\in\mathbb R$ for $B'\in\Mpsa$; continuity}
\lstep{2}{2}{$\ip{a-E'a}{i\Om}=i\bigl(\ip{a}{\Om}-\ip{E'a}{\Om}\bigr)$ and $\ip{E'a}{\Om}\in\mathbb R$.}
\lby{linearity in the second slot; $E'a\in\HMp$ by Lemma~\ref{lem:proj}, and $\ip{\zeta}{\Om}=
\overline{\ip{\Om}{\zeta}}\in\mathbb R$ for $\zeta\in\HMp$ by $\langle2\rangle1$}
\lstep{2}{3}{$\operatorname{Re}\ip{a-E'a}{i\Om}=\operatorname{Re}\bigl(i\ip{a}{\Om}\bigr)-0
=-\operatorname{Im}\ip{a}{\Om}=\operatorname{Im}\ip{\Om}{a}$.}
\lby{$\langle2\rangle2$ and $\operatorname{Re}(iz)=-\operatorname{Im}z$}
\lqed{2}{4}
\lby{$\langle2\rangle1$, $\langle2\rangle3$}
\lstep{1}{3}{$G_{\mathbb C}=G_{\mathbb R}-(\operatorname{Im}\ip{\Om}{a})^2$.}
\lby{by $\langle1\rangle2$ the closed real subspace $\HMp\oplus_{\mathbb R}\mathbb Ri\Om$ is an orthogonal direct sum,
$\norm{i\Om}=1$, and the residual $a-E'a$ has real-inner-product $\operatorname{Im}\ip{\Om}{a}$ with $i\Om$; so
projecting further onto $\mathbb Ri\Om$ removes exactly $(\operatorname{Im}\ip{\Om}{a})^2$ from the squared distance}
\lstep{1}{4}{The example $a=i\Om$.}
\lby{$\dist(i\Om,\HMp)^2=1$ because $\operatorname{Re}\ip{i\Om}{\zeta}=0$ for all $\zeta\in\HMp$ ($\langle2\rangle1$),
so $P_{\HMp}(i\Om)=0$; $G_{\mathbb C}=0$ by taking $c=i$; and
$\gB(i\Om)=2\norm{(1-J)u(i\Om)}^2=2\norm{iu\Om+iJu\Om}^2$ --- with $u\Om=(1+\Delta)^{-1/2}\Om=2^{-1/2}\Om$
(as $\Delta\Om=\Om$) and $J\Om=\Om$ --- $=2\norm{2\cdot i2^{-1/2}\Om}^2=4$; note $J(i2^{-1/2}\Om)=-i2^{-1/2}\Om$}
\lqed{1}{5}
\lby{$\langle1\rangle3$, $\langle1\rangle4$, Theorem~\ref{thm:three}, Lemma~\ref{lem:tangent}(3),(4)}

\section{(B) Finite dimensions}\label{sec:findim}

\begin{definition}[Finite-dimensional standard form]\label{def:findim}
Let $\rho=\sum_ip_i\ket i\bra i$ be a faithful density matrix on $\mathbb C^n$.  Realise the standard form on the
Hilbert--Schmidt space $\Hil=\mathfrak S_2(\mathbb C^n)$ of $n\times n$ matrices with $\ip XY=\Tr(X^*Y)$:
\begin{equation}\label{eq:findim-sf}
 \cM=\{L_X:X\in M_n\},\quad L_XY=XY;\qquad \cM'=\{R_X\},\quad R_XY=YX;\qquad
 \Om=\rho^{1/2};\qquad JY=Y^*;\qquad \Delta Y=\rho Y\rho^{-1}.
\end{equation}
This is the same standard form as Note~7's $\mathbb C^n\otimes\overline{\mathbb C^n}$ picture under
$\ket i\otimes\ket{\bar\jmath}\leftrightarrow\ket i\bra j$.
\end{definition}
\lstep{1}{1}{$\Om$ is cyclic and separating, $J$ is antiunitary with $J^2=1$, $J\Om=\Om$, $\Delta>0$, $\Delta\Om=\Om$.}
\lby{$\rho$ invertible, so $\cM\Om=M_n\rho^{1/2}=\Hil$ and $X\rho^{1/2}=0\Rightarrow X=0$;
$\Tr((Y^*)^*Z^*)=\Tr(YZ^*)=\overline{\Tr(ZY^*)}=\overline{\ip ZY}$, so $J$ is antiunitary;
$(\rho^{1/2})^*=\rho^{1/2}$; $\rho\rho^{1/2}\rho^{-1}=\rho^{1/2}$}
\lstep{1}{2}{$S=J\Delta^{1/2}$ satisfies $S(X\Om)=X^*\Om$; $\Delta$ has eigenvalue $p_i/p_j$ on $\ket i\bra j$;
$S^*=F=\Delta^{1/2}J$ acts as $FY=\rho^{1/2}Y^*\rho^{-1/2}$.}
\lby{$J\Delta^{1/2}(X\rho^{1/2})=J(\rho^{1/2}X\rho^{1/2}\rho^{-1/2})=J(\rho^{1/2}X)=X^*\rho^{1/2}$;
$\Delta(\ket i\bra j)=\rho\ket i\bra j\rho^{-1}=(p_i/p_j)\ket i\bra j$;
$\Delta^{1/2}JY=\rho^{1/2}Y^*\rho^{-1/2}$}

\begin{proposition}[Identification of the two metrics in finite dimensions]\label{prop:findim-id}
In the setting of Definition~\ref{def:findim} let $a\in\Hil$ with $\operatorname{Im}\ip\Om a=0$ and
$\varphi(X)=i\ip{a}{X\Om}-i\ip{X^*\Om}{a}$ for $X\in\cM$.  Let $d\rho$ be the unique matrix with
$\varphi(L_X)=\Tr(d\rho\,X)$ for all $X$.  Then
\begin{enumerate}[leftmargin=1.7em]
\item $d\rho=d\rho^*$, $\Tr d\rho=0$, $\eta:=i(F-1)a$ satisfies $\varphi(L_X)=\ip{\eta}{X\Om}$ and
 $\eta=d\rho\,\rho^{-1/2}$, i.e.\ $\eta_{ij}=d\rho_{ij}/\sqrt{p_j}$;
\item $\displaystyle \gB(a)=2\ip{\eta}{(1+\Delta)^{-1}\eta}=2\sum_{i,j}\frac{|d\rho_{ij}|^2}{p_i+p_j}$,
 the SLD quantum Fisher information;
\item $\displaystyle \gKM:=\Bigl\langle\eta,\frac{\log\Delta}{\Delta-1}\eta\Bigr\rangle
 =\sum_{i,j}|d\rho_{ij}|^2\,\ell(p_i,p_j)$, $\ \ell(p,q)=\frac{\log p-\log q}{p-q}$, $\ell(p,p)=1/p$;
\item $\tfrac14\gB(a)=\dist(a,\Mpsa\Om)^2$, the minimiser being $B'=R_Y$ with $\rho Y+Y\rho=\rho^{1/2}a+a^*\rho^{1/2}$,
 i.e.\ $Y=\tfrac12L$ with $L$ the symmetric logarithmic derivative $\rho L+L\rho=2\,d\rho$.
\end{enumerate}
\end{proposition}

\lstep{1}{1}{$\varphi(L_X)=-2\operatorname{Im}\Tr(a^*X\rho^{1/2})$, $\Tr d\rho=\varphi(\one)=0$, and $d\rho=d\rho^*$.}
\lby{$\ip{a}{X\Om}=\Tr(a^*X\rho^{1/2})$ and $\ip{X^*\Om}{a}=\Tr(\rho^{1/2}Xa)=\overline{\Tr(a^*X^*\rho^{1/2})}$;
hence for $X=X^*$, $\varphi=i(z-\bar z)=-2\operatorname{Im}z$; $\varphi(\one)=-2\operatorname{Im}\ip\Om a=0$;
$\varphi(L_X)$ is real for $X=X^*$, so $\Tr(d\rho X)\in\mathbb R$ for all Hermitian $X$, i.e.\ $d\rho$ Hermitian}
\lstep{1}{2}{Item (1).}
\lby{Lemma~\ref{lem:eta} (no domain condition in finite dimensions, $D(F)=\Hil$) gives
$\varphi(L_X)=\ip{\eta}{X\Om}$ with $\eta=i(F-1)a$; and $\ip{\eta}{X\Om}=\Tr(\eta^*X\rho^{1/2})
=\Tr(d\rho X)$ for all $X$ forces $\rho^{1/2}\eta^*=d\rho$, i.e.\ $\eta^*=\rho^{-1/2}d\rho$ and
$\eta=d\rho\,\rho^{-1/2}$ (using $d\rho=d\rho^*$)}
\lstep{1}{3}{Item (2).}
\lby{$(1+\Delta)^{-1}$ multiplies the matrix entry $(i,j)$ by $(1+p_i/p_j)^{-1}=p_j/(p_i+p_j)$, so
$\ip{\eta}{(1+\Delta)^{-1}\eta}=\sum_{ij}|\eta_{ij}|^2p_j/(p_i+p_j)=\sum_{ij}\frac{|d\rho_{ij}|^2}{p_j}\frac{p_j}{p_i+p_j}$;
and $\gB(a)=2\norm{(1+\Delta)^{-1/2}\eta}^2$ by Lemma~\ref{lem:eta}(3)}
\lstep{1}{4}{Item (3).}
\lby{$\frac{\log\Delta}{\Delta-1}$ multiplies the entry $(i,j)$ by
$\frac{\log(p_i/p_j)}{p_i/p_j-1}=p_j\frac{\log p_i-\log p_j}{p_i-p_j}=p_j\ell(p_i,p_j)$, and
$|\eta_{ij}|^2p_j=|d\rho_{ij}|^2$}
\lstep{1}{5}{Item (4).}
\lproof
\lstep{2}{1}{$\Mpsa\Om=\{\rho^{1/2}Y:Y=Y^*\}$, a closed real subspace.}
\lby{$R_Y$ is self-adjoint on $\Hil$ iff $Y=Y^*$ (as $\ip{R_YZ}{W}=\Tr(Y^*Z^*W)=\ip{Z}{R_{Y^*}W}$), and
$R_Y\Om=\rho^{1/2}Y$; the set is a finite-dimensional real subspace, hence closed}
\lstep{2}{2}{$h(Y):=\norm{a-\rho^{1/2}Y}^2=\norm a^2-\Tr\bigl(Y(\rho^{1/2}a+a^*\rho^{1/2})\bigr)+\Tr(Y\rho Y)$
for $Y=Y^*$.}
\lby{$\norm{a-\rho^{1/2}Y}^2=\Tr\bigl((a^*-Y\rho^{1/2})(a-\rho^{1/2}Y)\bigr)$ and cyclicity of the trace}
\lstep{2}{3}{$h$ is a strictly convex real quadratic form on the real vector space of Hermitian $Y$, and its unique
stationary point solves $\rho Y+Y\rho=R:=\rho^{1/2}a+a^*\rho^{1/2}$.}
\lby{the quadratic part $\Tr(Y\rho Y)=\norm{\rho^{1/2}Y}^2>0$ for $Y\ne0$ ($\rho$ invertible);
the derivative at $Y$ in a Hermitian direction $H$ is $-\Tr(HR)+\Tr(H\rho Y)+\Tr(Y\rho H)=\Tr\bigl(H(\rho Y+Y\rho-R)\bigr)$,
and $\Tr(HZ)=0$ for all Hermitian $H$ with $Z$ Hermitian forces $Z=0$}
\lstep{2}{4}{At the minimum, $h=\norm a^2-\tfrac12\Tr(YR)=\norm a^2-\tfrac12\sum_{ij}\frac{|R_{ij}|^2}{p_i+p_j}$
(entries in the eigenbasis of $\rho$).}
\lby{$\Tr(Y\rho Y)=\tfrac12\Tr(Y(\rho Y+Y\rho))=\tfrac12\Tr(YR)$ at stationarity; $Y_{ij}=R_{ij}/(p_i+p_j)$ solves
$\rho Y+Y\rho=R$ entrywise; $R$ Hermitian gives $\Tr(YR)=\sum_{ij}|R_{ij}|^2/(p_i+p_j)$}
\lstep{2}{5}{$\displaystyle\sum_{ij}\Bigl[|a_{ij}|^2-\tfrac12\frac{|R_{ij}|^2}{p_i+p_j}\Bigr]
 =\tfrac12\sum_{ij}\frac{|\eta_{ij}|^2p_j}{p_i+p_j}$, i.e.\ the minimum of $h$ equals $\tfrac14\gB(a)$.}
\lproof
\lstep{3}{1}{Put $\alpha=a_{ij}$, $\beta=\overline{a_{ji}}$, $\mu=\sqrt{p_i}$, $\nu=\sqrt{p_j}$.  Then
$R_{ij}=\mu\alpha+\nu\beta$ and $|\eta_{ij}|^2p_j=|\mu\beta-\nu\alpha|^2$.}
\lby{$R=\rho^{1/2}a+a^*\rho^{1/2}$ gives $R_{ij}=\sqrt{p_i}a_{ij}+\overline{a_{ji}}\sqrt{p_j}$;
$\eta=i(F-1)a$ with $(Fa)_{ij}=(\rho^{1/2}a^*\rho^{-1/2})_{ij}=\sqrt{p_i/p_j}\,\overline{a_{ji}}$, so
$\eta_{ij}=i(\sqrt{p_i/p_j}\overline{a_{ji}}-a_{ij})$ and $|\eta_{ij}|^2p_j=|\mu\beta-\nu\alpha|^2$}
\lstep{3}{2}{$|\mu\alpha+\nu\beta|^2+|\mu\beta-\nu\alpha|^2=(\mu^2+\nu^2)(|\alpha|^2+|\beta|^2)$.}
\lby{expanding, the cross terms are $2\mu\nu\operatorname{Re}(\bar\alpha\beta)-2\mu\nu\operatorname{Re}(\bar\beta\alpha)
=2\mu\nu\operatorname{Re}(\bar\alpha\beta-\overline{\bar\alpha\beta})=0$}
\lstep{3}{3}{The pair of index positions $(i,j)$ and $(j,i)$ contributes $|\alpha|^2+|\beta|^2$ to
$\sum|a_{ij}|^2$, $2|R_{ij}|^2/(p_i+p_j)$ (halved: $|R_{ij}|^2/(p_i+p_j)$ after the factor $\tfrac12$) to the second
sum, and $|\mu\beta-\nu\alpha|^2/(p_i+p_j)$ to the right-hand side.}
\lby{$|a_{ji}|^2=|\beta|^2$; $|R_{ji}|=|\overline{R_{ij}}|=|R_{ij}|$ since $R=R^*$;
$|\eta_{ji}|^2p_i=|\nu\bar\alpha-\mu\bar\beta|^2=|\mu\beta-\nu\alpha|^2$;
and for $i=j$ the same formulas hold with the pair counted once}
\lstep{3}{4}{The claim follows from $\langle3\rangle2$ divided by $\mu^2+\nu^2=p_i+p_j$.}
\lby{$|\alpha|^2+|\beta|^2-|R_{ij}|^2/(p_i+p_j)=|\mu\beta-\nu\alpha|^2/(p_i+p_j)$ by $\langle3\rangle1$, $\langle3\rangle2$}
\lqed{3}{5}
\lby{$\langle3\rangle3$, $\langle3\rangle4$, summed over all unordered pairs}
\lqed{2}{6}
\lby{$\langle2\rangle3$, $\langle2\rangle4$, $\langle2\rangle5$}
\lqed{1}{6}
\lby{$\langle1\rangle2$--$\langle1\rangle5$}

\begin{remark}[Two different stationarity equations, one value]\label{rem:SLD}
The minimiser of the \emph{distance} problem $\dist(a,\Mpsa\Om)$ solves $\rho Y+Y\rho=\rho^{1/2}a+a^*\rho^{1/2}$,
whereas the maximiser of the \emph{variational} problem $\sup_{K\in\Msa}[\varphi(K)-\Var_\omega(K)]$ solves the
symmetric-logarithmic-derivative equation $\rho K+K\rho=d\rho$, i.e.\ $K=\tfrac12L$ with $\rho L+L\rho=2\,d\rho$,
and has value $\tfrac12\Tr(d\rho\,K)=\tfrac14\Tr(d\rho\,L)=\tfrac14\sum_{ij}\frac{2|d\rho_{ij}|^2}{p_i+p_j}
=\tfrac14\gB$.  The two equations are different (the first involves $a$, the second the rotated tangent) but the two
values coincide, which is the content of Theorem~\ref{thm:three}.  The stationarity computation for the second is:
the derivative of $\Tr(d\rho K)-\Tr(\rho K^2)+\Tr(\rho K)^2$ in a Hermitian direction $H$ at a point with
$\omega(K)=0$ is $\Tr\bigl(H(d\rho-\rho K-K\rho)\bigr)$, which vanishes for all Hermitian $H$ iff
$d\rho=\rho K+K\rho$; and then $\Tr(\rho K^2)=\tfrac12\Tr(K(\rho K+K\rho))=\tfrac12\Tr(K\,d\rho)$, so the value is
$\Tr(d\rho K)-\tfrac12\Tr(d\rho K)=\tfrac12\Tr(d\rho K)$; finally $L_{ij}=2d\rho_{ij}/(p_i+p_j)$ solves the SLD
equation entrywise and $\Tr(d\rho L)=\sum_{ij}d\rho_{ij}L_{ji}=\sum_{ij}2|d\rho_{ij}|^2/(p_i+p_j)=\gB$
by Proposition~\ref{prop:findim-id}(2).
\end{remark}

\subsection{The Kubo--Mori second variation in finite dimensions}\label{sec:KM}

\begin{proposition}[Relative entropy]\label{prop:KM}
In the setting of Definition~\ref{def:findim} let $s\mapsto\rho_s$ be a $C^2$ curve of density matrices on
$\mathbb C^n$ with $\rho_0=\rho>0$, $\dot\rho_0=d\rho$.  Then
\begin{equation}\label{eq:KM-expansion}
 D(\rho\Vert\rho_s):=\Tr\bigl[\rho(\log\rho-\log\rho_s)\bigr]=\frac{s^2}{2}\gKM+O(s^3),\qquad
 \gKM=\sum_{i,j}|d\rho_{ij}|^2\ell(p_i,p_j)=\Bigl\langle\eta,\frac{\log\Delta}{\Delta-1}\eta\Bigr\rangle .
\end{equation}
In particular this applies to $\rho_s$ defined by $\Tr(\rho_sX)=\ip{\xi_s}{L_X\xi_s}$ with $\xi_s=e^{-isA}\Om$
and $A$ self-adjoint on $\Hil$, which is real-analytic in $s$.
\end{proposition}

\lstep{1}{1}{\ldefine $Y_s:=\rho_s-\rho=s\,d\rho+\tfrac{s^2}2d^2\rho+O(s^3)$ (norm), $R_t:=(\rho+t)^{-1}$ for $t\ge0$.
Then $\Tr Y_s=0$, $\Tr d\rho=0$, $\Tr d^2\rho=0$; and there are $s_1>0$, $C<\infty$ with
$\rho_s\ge\tfrac12p_{\min}$ and $\norm{Y_s}\le C|s|$ for $|s|\le s_1$, $p_{\min}=\min_ip_i>0$.}
\lby{Taylor's theorem with remainder for the $C^2$ curve $\rho_s$; $\Tr\rho_s=1$ for all $s$, so all $s$-derivatives of
$\Tr\rho_s$ vanish; continuity of the smallest eigenvalue}
\lstep{1}{2}{$\log X=\int_0^\infty\bigl[(1+t)^{-1}-(X+t)^{-1}\bigr]dt$ for $X>0$.}
\lby{spectral calculus and the elementary identity $\log x=\int_0^\infty\bigl[\frac1{1+t}-\frac1{x+t}\bigr]dt$ for
$x>0$, the integrand being $O(t^{-2})$ at infinity}
\lstep{1}{3}{$\displaystyle\Tr(\rho\log\rho_s)-\Tr(\rho\log\rho)
 =\int_0^\infty\Tr\bigl[\rho(R_t-(\rho_s+t)^{-1})\bigr]dt$.}
\lby{$\langle1\rangle2$ applied to $\rho_s$ and to $\rho$, subtracted; the $(1+t)^{-1}$ terms cancel; the $t$-integrals
converge absolutely because $\norm{R_t-(\rho_s+t)^{-1}}\le\norm{R_t}\norm{Y_s}\norm{(\rho_s+t)^{-1}}=O(t^{-2})$}
\lstep{1}{4}{$(\rho_s+t)^{-1}=R_t-R_tY_sR_t+R_tY_sR_tY_sR_t-\Theta_t$, with
$\norm{\Theta_t}\le\frac{\norm{R_t}^4\norm{Y_s}^3}{1-\norm{R_t}\norm{Y_s}}$ and
$\norm{R_t}\le(p_{\min}+t)^{-1}$.}
\lby{Neumann series for $(\rho+t+Y_s)^{-1}=\bigl(1+R_tY_s\bigr)^{-1}R_t$, convergent because
$\norm{R_tY_s}\le\norm{Y_s}/p_{\min}<1$ for $|s|$ small; the remainder after three terms is
$\sum_{k\ge3}(-1)^k(R_tY_s)^kR_t$}
\lstep{1}{5}{$\int_0^\infty\norm{\Theta_t}\,dt\le C'|s|^3$ for $|s|\le s_2$.}
\lby{$\int_0^\infty(p_{\min}+t)^{-4}dt=\tfrac13p_{\min}^{-3}<\infty$ and $\norm{Y_s}\le C|s|$ by $\langle1\rangle1$;
the denominator is $\ge\tfrac12$ for $|s|$ small}
\lstep{1}{6}{$\displaystyle\int_0^\infty\Tr(\rho R_tY_sR_t)\,dt=\Tr Y_s=0$.}
\lby{$\Tr(\rho R_tY_sR_t)=\Tr(Y_sR_t\rho R_t)$ and $\int_0^\infty R_t\rho R_t\,dt=\one$, because in the eigenbasis
$\int_0^\infty p_i(p_i+t)^{-2}dt=1$; then $\langle1\rangle1$}
\lstep{1}{7}{$\displaystyle\int_0^\infty\Tr(\rho R_tY_sR_tY_sR_t)\,dt=\frac{s^2}{2}\sum_{ij}|d\rho_{ij}|^2\ell(p_i,p_j)
 +O(s^3)$.}
\lproof
\lstep{2}{1}{$\Tr(\rho R_tY R_tYR_t)=\sum_{ij}|Y_{ij}|^2\dfrac{p_i}{(p_i+t)^2(p_j+t)}$ for Hermitian $Y$
(entries in the eigenbasis of $\rho$).}
\lby{$R_t$ is diagonal with entries $(p_i+t)^{-1}$; expanding the trace,
$\sum_{ij}p_i(p_i+t)^{-1}Y_{ij}(p_j+t)^{-1}Y_{ji}(p_i+t)^{-1}$, and $Y_{ji}=\overline{Y_{ij}}$}
\lstep{2}{2}{$\displaystyle\int_0^\infty\Bigl[\frac{p_i}{(p_i+t)^2(p_j+t)}+\frac{p_j}{(p_j+t)^2(p_i+t)}\Bigr]dt
 =\ell(p_i,p_j)$.}
\lby{the difference of the integrand and $\frac{1}{(p_i+t)(p_j+t)}$ is
$\frac{t^2-p_ip_j}{(p_i+t)^2(p_j+t)^2}=\frac{d}{dt}\Bigl[\frac{-t}{(p_i+t)(p_j+t)}\Bigr]$,
whose integral over $[0,\infty)$ is $0-0=0$; and $\int_0^\infty\frac{dt}{(p_i+t)(p_j+t)}
=\frac{\log p_i-\log p_j}{p_i-p_j}=\ell(p_i,p_j)$}
\lstep{2}{3}{$\displaystyle\int_0^\infty\Tr(\rho R_tYR_tYR_t)dt=\tfrac12\sum_{ij}|Y_{ij}|^2\ell(p_i,p_j)$.}
\lby{$\langle2\rangle1$, symmetrising the double sum in $i\leftrightarrow j$ (legitimate: $|Y_{ij}|=|Y_{ji}|$),
and $\langle2\rangle2$}
\lstep{2}{4}{With $Y=Y_s=s\,d\rho+O(s^2)$: $\sum_{ij}|(Y_s)_{ij}|^2\ell(p_i,p_j)=s^2\sum_{ij}|d\rho_{ij}|^2\ell(p_i,p_j)+O(s^3)$.}
\lby{$\ell(p_i,p_j)\le1/p_{\min}$ for all $i,j$, so the quadratic form is norm-continuous; expanding the square}
\lqed{2}{5}
\lby{$\langle2\rangle3$, $\langle2\rangle4$}
\lstep{1}{8}{$D(\rho\Vert\rho_s)=\frac{s^2}{2}\sum_{ij}|d\rho_{ij}|^2\ell(p_i,p_j)+O(s^3)$.}
\lby{$D=\Tr\rho\log\rho-\Tr\rho\log\rho_s=-\int\Tr[\rho(R_t-(\rho_s+t)^{-1})]dt$ by $\langle1\rangle3$;
insert $\langle1\rangle4$: $R_t-(\rho_s+t)^{-1}=R_tY_sR_t-R_tY_sR_tY_sR_t+\Theta_t$; integrate using
$\langle1\rangle6$ (first term, $=0$), $\langle1\rangle7$ (second term) and $\langle1\rangle5$ (remainder $O(s^3)$);
the overall sign: $D=-[0-\frac{s^2}2\sum|d\rho_{ij}|^2\ell+O(s^3)]$}
\lstep{1}{9}{$\sum_{ij}|d\rho_{ij}|^2\ell(p_i,p_j)=\ip{\eta}{\frac{\log\Delta}{\Delta-1}\eta}$.}
\lby{Proposition~\ref{prop:findim-id}(3)}
\lqed{1}{10}
\lby{$\langle1\rangle8$, $\langle1\rangle9$}

\begin{verdictok}
Checked numerically: for the random example of the verdict after Theorem~\ref{thm:three},
$D(\rho\Vert\rho_s)/s^2$ takes the values $17.513,\ 14.373,\ 14.029,\ 13.933$ at $s=10^{-1},3\cdot10^{-2},
10^{-2},3\cdot10^{-3}$, converging to $\gKM/2=13.895367$; the residual is $O(s)$ as predicted by the $O(s^3)$ error.
\end{verdictok}

\begin{remark}[Why the second-order term of $\rho_s$ drops out]
$\langle1\rangle6$ is the statement that the first-order derivative of $-\Tr\rho\log(\cdot)$ at $\rho$ is
$-\Tr(\cdot)$, so it annihilates \emph{every} traceless perturbation, in particular $\tfrac{s^2}2d^2\rho$.
This is why $D(\rho\Vert\rho_s)$ to order $s^2$ depends on the curve only through $d\rho$ --- exactly as
$1-\Fid$ does; it is also why no hypothesis such as $\Om\in D(A^2)$ enters the general statement.
\end{remark}

\section{(C) The Uhlmann--Alberti representations of the fidelity}\label{sec:uhlmann}

\begin{definition}[Transition probability]\label{def:P}
For a W$^*$-algebra $\cM$ and $\nu,\varrho\in\cM^*_+$, and writing $S_{\pi,\cM}(\nu)=\{\phi\in K:
\nu=\ip{\phi}{\pi(\cdot)\phi}\}$ for the set of vector representatives of $\nu$ in a unital $*$-representation
$\{\pi,K\}$,
\[
 P_{\cM}(\nu,\varrho)=\sup_{\{\pi,K\},\ \phi\in S_{\pi,\cM}(\nu),\ \psi\in S_{\pi,\cM}(\varrho)}|\ip{\psi}{\phi}|^2,
 \qquad \Fid(\nu,\varrho):=\sqrt{P_{\cM}(\nu,\varrho)} .
\]
In finite dimensions $\Fid(\rho,\sigma)=\Tr|\rho^{1/2}\sigma^{1/2}|$.
\end{definition}

\begin{citedbox}{Alberti--Uhlmann 2002, Definition 2 and estimates (1-4)--(1-6), p.~4}
\emph{Transcribed:} ``$P_{\cM}(\nu,\varrho)=\sup_{\{\pi,K\},\phi\in S_{\pi,\cM}(\nu),\psi\in S_{\pi,\cM}(\varrho)}
|\langle\psi,\phi\rangle|^2$'' (Definition~2); and
``$|f(\one)|^2\le P_{\cM}(\nu,\varrho)\le\nu(a)\varrho(a^{-1})$'' \ (1-4), where $f$ is any element of
$\Gamma_{\cM}(\nu,\varrho)=\{f\in\cM^*:|f(y^*x)|^2\le\nu(y^*y)\varrho(x^*x)\ \forall x,y\in\cM\}$ (1-5) and $a$ any
invertible positive element of $\cM_+$; and
``$\Gamma_{\cM}(\nu,\varrho)=\{\langle\pi(\cdot)k\psi,\phi\rangle:k\in(\pi(\cM)')_1\}$'' (1-6),
$(\pi(\cM)')_1$ the unit ball of the commutant.
\emph{Locator:} P.~Alberti, A.~Uhlmann, \emph{On Bures distance and $*$-algebraic transition probability between inner
derived positive linear forms}, arXiv:math-ph/0202038, p.~4 (journal: Rep.\ Math.\ Phys.\ 46 (2000) 97--104 is the
companion; this preprint is the extended version).  Local copy \texttt{refs/AlbertiUhlmann02\_math-ph-0202038.pdf}.
\emph{Screenshot:}\\
\includegraphics[width=\textwidth]{shots/P1_AU02_def2.png}\\
\includegraphics[width=\textwidth]{shots/P1_AU02_est.png}\\
\emph{Use here:} the upper estimate of (1-4) is the \emph{only} input for the upper bound (E),
Section~\ref{sec:upper}; the lower estimate of (1-4) together with (1-6) (take $k=v'$ a unitary of $\cM'$) is the
only input for the lower bound (D), Section~\ref{sec:lower}.  \emph{Verdict on usage:} both are applied for normal
\emph{states} of a von Neumann algebra in standard form, a special case of the quoted general setting; the
inner-product convention of \cite{AU02} is linear in the first variable, which is irrelevant since only
$|\ip\psi\phi|$ enters.
\end{citedbox}

\begin{citedbox}{Alberti--Uhlmann 2002, Theorem 2, p.~6}
\emph{Transcribed:} ``THEOREM 2. Let $M$ be a W$^*$-algebra, and be $\nu,\varrho\in M^*_+$. Then, the following facts
hold: (1) $\sqrt{P_M(\nu,\varrho)}=\inf_{x>0}\sqrt{\nu(x)\varrho(x^{-1})}$; (2)
$\sqrt{P_M(\nu,\varrho)}=\sup_{f\in\Gamma_M(\nu,\varrho)}|f(\one)|$.  The infimum in (1) extends over all positive
invertible elements of $M$.  Moreover, if $\{\pi,K\}$ is any unital $*$-representation of $M$ over some Hilbert space
$K$ such that $S_{\pi,M}(\nu)\ne\emptyset$ and $S_{\pi,M}(\varrho)\ne\emptyset$ are fulfilled, then the following is
fulfilled: (3) $\sqrt{P_M(\nu,\varrho)}=\sup_{\psi\in S_{\pi,M}(\varrho)}|\langle\psi,\phi\rangle|$,
$\forall\phi\in S_{\pi,M}(\nu)$.  Also, the supremum in (2) is a maximum\dots The previous result remains valid even
if $M$ is supposed to be a unital C$^*$-algebra.''
\emph{Locator:} arXiv:math-ph/0202038, p.~6; proofs are attributed there to Alberti 1983 (Corollary 1, Corollary 3,
Theorem 3 of that paper) --- Alberti 1983 itself is paywalled and \textbf{NOT OBTAINED}; this open preprint is the
restatement we rely on.
\emph{Screenshot:}\\
\includegraphics[width=\textwidth]{shots/P1_AU02_thm2.png}\\
\emph{Use here:} Corollary~\ref{cor:uhlmann} below.  \emph{Verdict:} hypotheses are satisfied ($\cM$ a von Neumann
algebra, $\nu=\omega$ and $\varrho=\omega_s$ normal states, $\pi=\mathrm{id}$ on the standard-form space, both
representative sets nonempty).
\end{citedbox}

\begin{corollary}[Uhlmann representation in standard form]\label{cor:uhlmann}
Let $\omega=\ip\Om{\cdot\,\Om}$ be faithful normal on $\cM$ with $\Om$ cyclic and separating, and let $\varrho$ be a
normal state of $\cM$ with a vector representative $\xi\in\Hil$.  Then
\begin{equation}\label{eq:uhlmann}
 \Fid(\omega,\varrho)=\sup\bigl\{|\ip{\Om}{v'\xi}|:\ v'\in\U(\cM')\bigr\},
\end{equation}
$\U(\cM')$ the unitary group of $\cM'$.  Moreover
\begin{equation}\label{eq:bures-dist}
 2\bigl(1-\Fid(\omega,\varrho)\bigr)=\inf\bigl\{\norm{\Om-v'\xi}^2:\ v'\in\U(\cM')\bigr\}.
\end{equation}
\end{corollary}
\lstep{1}{1}{$\{\phi\in\Hil:\ \omega=\ip\phi{\cdot\,\phi}\}=\U(\cM')\Om$.}
\lproof
\lstep{2}{1}{$\supseteq$: $\ip{v'\Om}{Xv'\Om}=\ip{\Om}{v'^*Xv'\Om}=\ip\Om{X\Om}$ for $X\in\cM$, $v'\in\U(\cM')$.}
\lby{$v'$ commutes with $X$ and is unitary}
\lstep{2}{2}{$\subseteq$: if $\omega_\phi=\omega_\Om$ on $\cM$ there is a partial isometry $v'\in\cM'$ with
$v'\Om=\phi$, $v'^*v'=[\cM'\Om]$, $v'v'^*=[\cM'\phi]$.}
\lby{Bratteli--Robinson, Prop.~2.5.30 (uniqueness of vector representatives up to a partial isometry in the commutant);
equivalently Dixmier, \emph{von Neumann Algebras}, Ch.~I \S4 Thm.~3}
\lstep{2}{3}{$[\cM'\Om]=\one$ and $[\cM'\phi]=\one$, so $v'$ is unitary.}
\lby{$\Om$ is cyclic for $\cM'$ because it is separating for $\cM$; $\phi$ is separating for $\cM$ because
$\omega_\phi=\omega$ is faithful ($\norm{X\phi}^2=\omega(X^*X)=0\Rightarrow X=0$), hence cyclic for $\cM'$
(Bratteli--Robinson Prop.~2.5.3: $\phi$ cyclic for $\cM'$ $\Leftrightarrow$ $\phi$ separating for $\cM$)}
\lqed{2}{4}
\lby{$\langle2\rangle1$--$\langle2\rangle3$}
\lstep{1}{2}{$\Fid(\omega,\varrho)=\sup\{|\ip{\xi}{\phi}|:\phi\in S_{\mathrm{id},\cM}(\omega)\}
 =\sup\{|\ip{\xi}{v'\Om}|:v'\in\U(\cM')\}$.}
\lby{Alberti--Uhlmann Theorem 2(3) with the roles $\nu=\varrho$, $\varrho=\omega$, the fixed representative
$\psi:=\xi\in S_{\mathrm{id},\cM}(\varrho)\ne\emptyset$ and the supremum running over $S_{\mathrm{id},\cM}(\omega)$;
then $\langle1\rangle1$}
\lstep{1}{3}{\eqref{eq:uhlmann}.}
\lby{$|\ip{\xi}{v'\Om}|=|\overline{\ip{v'\Om}{\xi}}|=|\ip{\Om}{v'^*\xi}|$ and $v'\mapsto v'^*$ is a bijection of
$\U(\cM')$}
\lstep{1}{4}{\eqref{eq:bures-dist}.}
\lby{for unit vectors, $\norm{\Om-v'\xi}^2=2-2\operatorname{Re}\ip{\Om}{v'\xi}$; and
$\sup_{v'}\operatorname{Re}\ip{\Om}{v'\xi}=\sup_{v'}|\ip{\Om}{v'\xi}|$ because $e^{i\theta}v'\in\U(\cM')$ for every
$\theta$, so the phase can be rotated; then $\langle1\rangle3$}
\lqed{1}{5}
\lby{$\langle1\rangle3$, $\langle1\rangle4$}

\begin{issue}[The proof route quoted in Note~7 is not available]\label{iss:wrong-uhlmann}
Note~7 writes, immediately after its Lemma~3.2: ``\emph{Proof route: Uhlmann's transition probability in standard
form, $\Fid(\omega,\omega_s)=\langle\Om,\Delta_{\omega_s,\omega}^{1/2}\Om\rangle$, and Araki's perturbation theory of
relative modular operators; alternatively Jen\v cov\'a's construction of monotone metrics through the standard
purification.}''  Three objections.
\begin{enumerate}[leftmargin=1.7em]
\item The displayed formula is \textbf{false}.  In finite dimensions
$\ip{\Om}{\Delta_{\sigma,\rho}^{1/2}\Om}=\Tr(\rho^{1/2}\sigma^{1/2})$, the Wigner--Yanase / $Q_{1/2}$ quasi-entropy,
whereas $\Fid=\Tr|\rho^{1/2}\sigma^{1/2}|=\Tr\bigl[(\rho^{1/2}\sigma\rho^{1/2})^{1/2}\bigr]$.  These agree iff
$\rho^{1/2}\sigma^{1/2}\ge0$; in general $\Tr(\rho^{1/2}\sigma^{1/2})<\Fid$, and the second variation of the former is
the Wigner--Yanase metric $g_{\mathrm{WY}}=4\sum_{ij}|d\rho_{ij}|^2/(\sqrt{p_i}+\sqrt{p_j})^2$, \emph{not} $\gB$.
(Independently established by agent~R3 of this review.)
\item Consequently the natural-cone representatives do \emph{not} compute the fidelity: for $\xi_\rho,\xi_\sigma$
in the natural cone $\mathcal P^\natural$ one has $\ip{\xi_\rho}{\xi_\sigma}=\Tr(\rho^{1/2}\sigma^{1/2})\le\Fid$, and
the supremum in \eqref{eq:uhlmann} is in general \emph{not} attained at $v'=\one$.  Any argument that replaces the
supremum over $\U(\cM')$ by the natural-cone pair is wrong.
\item Jen\v cov\'a, \emph{Quantum information geometry and standard purification}, JMP 43 (2002) 2187 works with
finite-dimensional (matrix) state spaces in the parts relevant here, so it cannot be cited for a type-III statement.
\end{enumerate}
The lemma is nevertheless true for the fidelity; Sections~\ref{sec:lower}--\ref{sec:upper} prove it by the route
$\Fid=\sup_{\U(\cM')}=\inf_{x>0}$ of Corollary~\ref{cor:uhlmann} and Alberti--Uhlmann Theorem~2(1).
\end{issue}

\section{(D) The lower bound on the fidelity}\label{sec:lower}

\begin{theorem}[Lower bound]\label{thm:lower}
Assume hypothesis \textup{(V)} with tangent representative $a=i\dot\xi_0$.  Then
\begin{equation}\label{eq:lower}
 \limsup_{s\to0}\frac{1-\Fid(\omega,\omega_s)}{s^2}\ \le\ \frac18\gB(a),
 \qquad\text{i.e.}\qquad \Fid(\omega,\omega_s)\ \ge\ 1-\frac{s^2}{8}\gB(a)+o(s^2).
\end{equation}
\end{theorem}

\lstep{1}{1}{\lpick $\eps>0$ and $B'\in\Mpsa$ with $\norm{a+B'\Om}^2\le\dist(a,\HMp)^2+\eps$.}
\lby{$\HMp=\overline{\Mpsa\Om}$ and $-B'$ runs through $\Mpsa$ when $B'$ does, so
$\inf_{B'\in\Mpsa}\norm{a+B'\Om}=\dist(a,\HMp)$, and an infimum is approached}
\lstep{1}{2}{$v'_s:=e^{-isB'}\in\U(\cM')$ and $\norm{e^{isB'}\Om-\Om-isB'\Om}\le\tfrac{s^2}2\norm{B'}^2$.}
\lby{$B'\in\cM'$ bounded self-adjoint, so $e^{-isB'}$ is a unitary of $\cM'$;
$\norm{(e^{iz}-1-iz)}\le z^2/2$ for real $z$ gives, by the spectral theorem for $B'$,
$\norm{(e^{isB'}-1-isB')\Om}\le\tfrac{s^2}2\norm{B'}^2\norm\Om$}
\lstep{1}{3}{$\norm{\Om-v'_s\xi_s}=\norm{e^{isB'}\Om-\xi_s}=|s|\,\norm{a+B'\Om}+o(s)$.}
\lproof
\lstep{2}{1}{$\norm{\Om-v'_s\xi_s}=\norm{v'^*_s(\Om-v'_s\xi_s)}=\norm{e^{isB'}\Om-\xi_s}$.}
\lby{$v'^*_s$ unitary}
\lstep{2}{2}{$e^{isB'}\Om-\xi_s=is\bigl(B'\Om+a\bigr)+o(s)$.}
\lby{$e^{isB'}\Om=\Om+isB'\Om+O(s^2)$ by $\langle1\rangle2$, and $\xi_s=\Om+s\dot\xi_0+o(s)=\Om-isa+o(s)$ by
hypothesis (V) and $a=i\dot\xi_0$, i.e.\ $\dot\xi_0=-ia$; subtracting,
$e^{isB'}\Om-\xi_s=isB'\Om+isa+o(s)$}
\lstep{2}{3}{$\norm{is(B'\Om+a)+o(s)}=|s|\norm{B'\Om+a}+o(s)$.}
\lby{the reverse triangle inequality $\bigl|\norm{\alpha+\beta}-\norm\alpha\bigr|\le\norm\beta$ with
$\alpha=is(B'\Om+a)$, $\norm\alpha=|s|\norm{B'\Om+a}$}
\lqed{2}{4}
\lby{$\langle2\rangle1$--$\langle2\rangle3$}
\lstep{1}{4}{$2(1-\Fid(\omega,\omega_s))\le\norm{\Om-v'_s\xi_s}^2$.}
\lby{Corollary~\ref{cor:uhlmann}, \eqref{eq:bures-dist}, whose right-hand side is an infimum over $\U(\cM')$;
$\xi_s$ is a vector representative of $\omega_s$ by hypothesis (V), and $\omega$ is faithful normal with cyclic
separating $\Om$}
\lstep{1}{5}{$\displaystyle\limsup_{s\to0}\frac{1-\Fid(\omega,\omega_s)}{s^2}\le\tfrac12\bigl(\dist(a,\HMp)^2+\eps\bigr)$.}
\lby{$\langle1\rangle3$, $\langle1\rangle4$: $2(1-\Fid)\le s^2\norm{a+B'\Om}^2+o(s^2)$, divide by $2s^2$ and use
$\langle1\rangle1$}
\lstep{1}{6}{\eqref{eq:lower}.}
\lby{$\eps>0$ arbitrary in $\langle1\rangle5$, and $\dist(a,\HMp)^2=\tfrac14\gB(a)$ by
Lemma~\ref{lem:proj}/Theorem~\ref{thm:three}}
\lqed{1}{7}
\lby{$\langle1\rangle5$, $\langle1\rangle6$}

\begin{remark}[No phase problem, no $\Om\in D(A^2)$]
The task's formulation of (D) asks to control $\ip{\Om}{e^{isB'}e^{-isA}\Om}$ to second order and to handle the
first-order imaginary term.  That is unnecessary: \eqref{eq:bures-dist} has already replaced $|\ip{\Om}{v'\xi_s}|$ by
a \emph{distance}, in which the phase is absorbed into $v'$ (step $\langle1\rangle4$ of Corollary~\ref{cor:uhlmann}),
and only the \emph{first} derivative of the curve is then needed.  In particular no hypothesis such as $\Om\in D(A^2)$,
and no product formula for $e^{isB'}e^{-isA}$ (which would require domain conditions on $A+B'$), enters.
\end{remark}

\section{(E) The upper bound on the fidelity}\label{sec:upper}

\begin{theorem}[Upper bound]\label{thm:upper}
Assume hypothesis \textup{(F$\varphi$)} with tangent representative $a$ \textup{(}so
$\varphi(K)=-2\operatorname{Im}\ip a{K\Om}$ for $K\in\Msa$ and $\operatorname{Im}\ip\Om a=0$\textup{)}.  Then
\begin{equation}\label{eq:upper}
 \liminf_{s\to0}\frac{1-\Fid(\omega,\omega_s)}{s^2}\ \ge\ \frac18\gB(a),
 \qquad\text{i.e.}\qquad \Fid(\omega,\omega_s)\ \le\ 1-\frac{s^2}{8}\gB(a)+o(s^2).
\end{equation}
\end{theorem}

\lstep{1}{1}{\lpick $K\in\Msa$ and let $|s|<\norm K^{-1}$.  Then $x_s:=\one+sK\in\cM$ is positive and invertible,
$x_s^{-1}=\one-sK+s^2K^2+r_s$ with $\norm{r_s}\le\dfrac{|s|^3\norm K^3}{1-|s|\,\norm K}$.}
\lby{$\norm{sK}<1$, so $x_s\ge(1-|s|\norm K)\one>0$ and the Neumann series
$x_s^{-1}=\sum_{k\ge0}(-sK)^k$ converges in norm; $r_s=\sum_{k\ge3}(-sK)^k$}
\lstep{1}{2}{$\Fid(\omega,\omega_s)^2\le\omega(x_s)\,\omega_s(x_s^{-1})$.}
\lby{Alberti--Uhlmann (1-4), upper estimate: $P_{\cM}(\nu,\varrho)\le\nu(x)\varrho(x^{-1})$ for every positive
invertible $x\in\cM$; here $\nu=\omega$, $\varrho=\omega_s$, $x=x_s$, and $\Fid^2=P$ by Definition~\ref{def:P}
(alternatively Theorem~2(1) of the same paper, of which this is the ``$\le$'' half)}
\lstep{1}{3}{$\omega_s(x_s^{-1})=1-s\,\omega(K)-s^2\varphi(K)+s^2\omega(K^2)+o(s^2)$.}
\lproof
\lstep{2}{1}{$\omega_s(x_s^{-1})=1-s\,\omega_s(K)+s^2\omega_s(K^2)+\omega_s(r_s)$ and $|\omega_s(r_s)|\le\norm{r_s}=O(s^3)$.}
\lby{$\langle1\rangle1$, linearity of $\omega_s$, $\omega_s(\one)=1$, and $|\omega_s(y)|\le\norm y$ for a state}
\lstep{2}{2}{$s\,\omega_s(K)=s\,\omega(K)+s^2\varphi(K)+o(s^2)$ and $s^2\omega_s(K^2)=s^2\omega(K^2)+o(s^2)$.}
\lby{hypothesis (F$\varphi$) applied to $y=K$ gives $\omega_s(K)=\omega(K)+s\varphi(K)+o(s)$; applied to $y=K^2$ it
gives in particular $\omega_s(K^2)\to\omega(K^2)$}
\lqed{2}{3}
\lby{$\langle2\rangle1$, $\langle2\rangle2$}
\lstep{1}{4}{$\omega(x_s)\,\omega_s(x_s^{-1})=1-s^2\bigl[\varphi(K)-\Var_\omega(K)\bigr]+o(s^2)$.}
\lby{$\omega(x_s)=1+s\omega(K)$ exactly; multiplying with $\langle1\rangle3$,
the $O(s)$ terms cancel ($s\omega(K)-s\omega(K)$) and the $s^2$ terms are
$-\varphi(K)+\omega(K^2)-\omega(K)^2=-[\varphi(K)-\Var_\omega(K)]$}
\lstep{1}{5}{$1-\Fid(\omega,\omega_s)\ \ge\ \tfrac12\bigl(1-\Fid(\omega,\omega_s)^2\bigr)
 \ \ge\ \tfrac{s^2}2\bigl[\varphi(K)-\Var_\omega(K)\bigr]+o(s^2)$.}
\lby{$1-\Fid^2=(1-\Fid)(1+\Fid)\le2(1-\Fid)$ because $0\le\Fid\le1$; then $\langle1\rangle2$, $\langle1\rangle4$}
\lstep{1}{6}{$\displaystyle\liminf_{s\to0}\frac{1-\Fid(\omega,\omega_s)}{s^2}\ge
 \tfrac12\bigl[\varphi(K)-\Var_\omega(K)\bigr]$ for every $K\in\Msa$.}
\lby{$\langle1\rangle5$, divided by $s^2>0$; the $o(s^2)$ depends on $K$, but $K$ is fixed while $s\to0$}
\lqed{1}{7}
\lby{taking the supremum over $K\in\Msa$ in $\langle1\rangle6$ and using
$\sup_{K\in\Msa}[\varphi(K)-\Var_\omega(K)]=\tfrac14\gB(a)$, which is Theorem~\ref{thm:three}}

\begin{theorem}[Type-III second variation of the fidelity: Note~7 Lemma~3.2, fidelity half]\label{thm:main}
Let $\cM$ be a von Neumann algebra with a cyclic and separating vector $\Om$ and $\omega=\ip\Om{\cdot\,\Om}$.
Let $(\omega_s)$ satisfy hypothesis \textup{(V)} \textup{(}which implies \textup{(F$\varphi$)}\textup{)}, with tangent
representative $a\in\Hil$; e.g.\ $\omega_s=\ip{u_s^*\Om}{\cdot\,u_s^*\Om}$ with $u_s=e^{isA}$ a strongly continuous
unitary group, $A=A^*$ on $\Hil$, $\Om\in D(A)$, $a=A\Om$.  Then
\begin{equation}\label{eq:main}
 \boxed{\ \Fid(\omega,\omega_s)=1-\frac{s^2}{8}\gB+o(s^2),\qquad
 -\log\Fid(\omega,\omega_s)=\frac{s^2}{8}\gB+o(s^2),\ }
\end{equation}
with
\begin{equation}\label{eq:main-gB}
 \gB=2\bigl\|(1-J)(1+\Delta)^{-1/2}a\bigr\|^2
   =4\dist\bigl(a,\overline{\Mpsa\Om}\bigr)^2
   =4\sup_{K\in\Msa}\bigl[\varphi(K)-\Var_\omega(K)\bigr]
   \;\bigl(=2\ip{\eta}{(1+\Delta)^{-1}\eta}\ \text{if }a\in D(S^*)\bigr).
\end{equation}
No hypothesis on the type of $\cM$, no separability, no hyperfiniteness and no assumption beyond $\Om\in D(A)$
are used.
\end{theorem}
\lstep{1}{1}{$\lim_{s\to0}s^{-2}(1-\Fid)=\tfrac18\gB$.}
\lby{Theorem~\ref{thm:lower} ($\limsup\le$) and Theorem~\ref{thm:upper} ($\liminf\ge$); (V) $\Rightarrow$ (F$\varphi$)
by Lemma~\ref{lem:tangent}(4), with the same $a$ and $\operatorname{Im}\ip\Om a=0$}
\lstep{1}{2}{$-\log\Fid=(1-\Fid)+O\bigl((1-\Fid)^2\bigr)=\tfrac{s^2}8\gB+o(s^2)$.}
\lby{$-\log(1-x)=x+O(x^2)$ for $x\to0$, with $x=1-\Fid\to0$ by $\langle1\rangle1$; and $O(s^4)=o(s^2)$}
\lstep{1}{3}{\eqref{eq:main-gB}.}
\lby{Definition~\ref{def:gB}, Theorem~\ref{thm:three}, Lemma~\ref{lem:eta}(3)}
\lqed{1}{4}
\lby{$\langle1\rangle1$--$\langle1\rangle3$}

\begin{verdictok}
Note~7's Lemma~3.2 is \emph{true for the fidelity} and is proved here in a stronger form than stated there
(hypothesis (V) rather than a unitary group; no condition on $\eta$; no hyperfiniteness).  Its stated
\emph{proof route} is not usable, see Issue~\ref{iss:wrong-uhlmann}.  Note~7's normalisation
$-\log\Fid=\frac{s^2}8\gB+o(s^2)$ and its formula $\gB=2\norm{(1+\Delta)^{-1/2}(S^*-1)A\Om}^2
=2\norm{(1-J)(1+\Delta)^{-1/2}A\Om}^2$ are confirmed, provided $A\Om\in\Hil$
(see Section~\ref{sec:open}, Gap~2, for the zero-collar case where this fails).
\end{verdictok}

\section{Monotonicity and the finite-dimensional approximation}\label{sec:mono}

The route proposed in item (E)(i) of the task --- exhaust a hyperfinite $\cM$ by finite-dimensional $\cM_n$, use
$\Fid_{\cM_n}\downarrow\Fid_{\cM}$ and the finite-dimensional expansion on each $\cM_n$ --- is not needed after
Theorem~\ref{thm:upper}.  It is however worth recording what it does and does not give.

\begin{proposition}[Data processing and monotone convergence of $\gB$]\label{prop:mono}
Let $\cM_1\subseteq\cM_2\subseteq\cdots\subseteq\cM$ be von Neumann subalgebras containing $\one$ with
$\bigl(\bigcup_n\cM_n\bigr)''=\cM$, let $\omega,\varphi,a$ be as in Theorem~\ref{thm:three}, and let
$\gB^{(n)}:=4\sup\{\varphi(K)-\Var_\omega(K):K\in(\cM_n)_{\mathrm{sa}}\}$ be the Bures form of the restricted data
$(\cM_n,\omega|_{\cM_n},\varphi|_{\cM_n})$.  Then
\begin{equation}\label{eq:mono}
 \gB^{(1)}\le\gB^{(2)}\le\cdots\le\gB,\qquad \lim_n\gB^{(n)}=\gB .
\end{equation}
\end{proposition}
\lstep{1}{1}{$\gB^{(n)}$ is well defined by the displayed variational expression, and coincides with
$2\norm{(1-J_n)(1+\Delta_n)^{-1/2}a}^2$ whenever $\Om$ is cyclic and separating for $\cM_n$, $(\Delta_n,J_n)$
being the modular data of $(\cM_n,\Om)$.}
\lby{the variational expression involves only $\varphi|_{(\cM_n)_{\mathrm{sa}}}$ and $\omega|_{\cM_n}$, hence makes
sense for an arbitrary unital subalgebra; the identification is Theorem~\ref{thm:three} applied to $(\cM_n,\Om)$,
whose hypotheses ($\Om$ cyclic and separating for $\cM_n$, $\operatorname{Im}\ip\Om a=0$) must be checked separately
--- they are \emph{not} automatic for a subalgebra, which is why the variational expression is taken as the definition}

\lstep{1}{2}{Monotonicity.}
\lby{$(\cM_n)_{\mathrm{sa}}\subseteq(\cM_{n+1})_{\mathrm{sa}}\subseteq\Msa$, and a supremum over a larger set is larger}
\lstep{1}{3}{$\sup_n\gB^{(n)}\ge\gB$.}
\lproof
\lstep{2}{1}{\lpick $K\in\Msa$ and $\delta>0$.  There is a net $(K_j)$ in $\bigcup_n(\cM_n)_{\mathrm{sa}}$ with
$\norm{K_j}\le\norm K$ and $K_j\to K$ in the strong$^*$ topology.}
\lby{Kaplansky's density theorem in the form: the self-adjoint part of the unit ball of a $*$-subalgebra
$\mathcal A$ is strongly dense in the self-adjoint part of the unit ball of $\mathcal A''$
(Bratteli--Robinson Thm.~2.4.16); apply it to $\mathcal A=\bigcup_n\cM_n$, which is a $*$-algebra because the
$\cM_n$ increase; on bounded sets strong and strong$^*$ convergence coincide for self-adjoint elements}
\lstep{2}{2}{$\varphi(K_j)\to\varphi(K)$ and $\Var_\omega(K_j)\to\Var_\omega(K)$.}
\lby{$K_j\Om\to K\Om$ (strong convergence) gives $\varphi(K_j)=-2\operatorname{Im}\ip a{K_j\Om}\to\varphi(K)$,
$\omega(K_j)=\ip\Om{K_j\Om}\to\omega(K)$ and $\omega(K_j^2)=\norm{K_j\Om}^2\to\norm{K\Om}^2=\omega(K^2)$}
\lstep{2}{3}{There is $j$ with $\varphi(K_j)-\Var_\omega(K_j)\ge\varphi(K)-\Var_\omega(K)-\delta$, and
$K_j\in(\cM_{n})_{\mathrm{sa}}$ for some $n$.}
\lby{$\langle2\rangle2$ and the definition of a net limit; $\bigcup_n\cM_n$ is the union}
\lqed{2}{4}
\lby{$\langle2\rangle3$ gives $\sup_n\gB^{(n)}\ge4[\varphi(K)-\Var_\omega(K)]-4\delta$; take the supremum over $K$
and let $\delta\downarrow0$}
\lqed{1}{4}
\lby{$\langle1\rangle2$, $\langle1\rangle3$}

\begin{remark}[Data processing]
Proposition~\ref{prop:mono} is the data-processing inequality for the SLD/Bures metric under the restriction
channel $\cM\to\cM_n$, and its proof through the variational expression of Theorem~\ref{thm:three} is one line:
\emph{the supremum is over a smaller set}.  The same argument proves data processing under an arbitrary unital normal
$2$-positive $\Phi:\cN\to\cM$ once one knows $\Var_{\omega\circ\Phi}(\Phi(K))\le\Phi(\Var)$ in the Kadison--Schwarz
form $\Phi(K)^2\le\Phi(K^2)$, i.e.\ $\Var_\omega(\Phi(K))\le\Var_{\omega\circ\Phi}(K)$; we do not need this here.
\end{remark}

\begin{remark}[Where the hyperfinite route stalls]\label{rem:hyperfinite}
Suppose $\cM$ is hyperfinite (as are all local algebras of a conformal net: they are the unique hyperfinite
type~III$_1$ factor --- see Section~\ref{sec:note7}) and $\cM_n\uparrow\cM$ with $\cM_n$ finite-dimensional.
By Note~4, Lemma~10.1 (fidelity is continuous from below), $\Fid_{\cM_n}(\omega,\omega_s)\downarrow\Fid_{\cM}(\omega,\omega_s)$
for each fixed $s$, and by the finite-dimensional case (Theorem~\ref{thm:main} applied to $\cM_n$),
$1-\Fid_{\cM_n}=\frac{s^2}8\gB^{(n)}+o_n(s^2)$.  To conclude
$1-\Fid_{\cM}=\frac{s^2}8\gB+o(s^2)$ one must interchange $\lim_n$ and $\lim_{s\to0}$, i.e.\ one needs the remainders
$o_n(s^2)$ to be \emph{uniform in $n$}.  They are not, in any evident way: the finite-dimensional remainder is
controlled by quantities such as $\norm{L_n}$ (the SLD of the compressed data) and $p^{(n)}_{\min}$, both of which
blow up as $n\to\infty$ for a type~III algebra --- indeed $p^{(n)}_{\min}\to0$ necessarily, since $\cM$ has no
normal trace.  Moreover the inequality goes the wrong way for the direction one wants:
$\Fid_{\cM_n}\ge\Fid_{\cM}$ gives $1-\Fid_{\cM}\ge1-\Fid_{\cM_n}=\frac{s^2}8\gB^{(n)}+o_n(s^2)$, which yields
$\liminf_s s^{-2}(1-\Fid_{\cM})\ge\frac18\gB^{(n)}$ for each $n$ and hence, by Proposition~\ref{prop:mono},
$\liminf_s s^{-2}(1-\Fid_{\cM})\ge\frac18\gB$ --- \emph{this} direction (the upper bound on $\Fid$) does close, and is
an independent proof of Theorem~\ref{thm:upper} in the hyperfinite case.  The other direction, the lower bound on
$\Fid$, does \emph{not} follow from the approximation and needs Theorem~\ref{thm:lower} anyway.
So the hyperfinite route gives at best half of the lemma, and exactly the half that Alberti's inequality gives in
complete generality.
\end{remark}

\begin{verdictok}
Note~4, Lemma~10.1 (``$\Fid_{M_n}(\varphi|_{M_n},\psi|_{M_n})\downarrow\Fid_M(\varphi,\psi)$'') is correct.  Its proof
uses the variational formula $\Fid_M(\varphi,\psi)=\inf_{x\in M,x>0}\frac12(\varphi(x)+\psi(x^{-1}))$, which is
Alberti--Uhlmann Corollary~2(5) (p.~8 of the preprint), and the reduction to $m\one\le x\le M_0\one$ is automatic,
since a positive invertible $x$ of a C$^*$-algebra satisfies $\norm{x^{-1}}^{-1}\one\le x\le\norm x\one$.
The Kaplansky-density step is as in $\langle2\rangle1$ above, and $x_k\to x$, $x_k^{-1}\to x^{-1}$ strongly with
uniform bounds is standard (continuity of the inverse on the strong topology restricted to a set with a uniform
lower bound).
\end{verdictok}

\section{(F) Application to the curve of Note~7}\label{sec:note7}

\subsection{The setting of Note~7}
Note~7 (Section~1, ``The deformation'') considers a diffeomorphism covariant chiral net $\cA$ on $S^1=\mathbb R\cup\{\infty\}$,
the intervals $A=(-\infty,0)$, $A_\eps=(-\infty,-\eps)$, $D=(0,L)$, $I=(-\infty,L)$, the M\"obius flow
$h_s(x)=Lx/(L+sx)$ with generator $f(x)=-x^2/L$, $U(s)=V(h_s)=e^{isT(f)}$, and the deformed state
\begin{equation}\label{eq:N7-state}
 \widetilde\omega_\eps(x_Ay_D)=\omega\bigl(x_A\,U(s)y_DU(s)^*\bigr),\qquad x_A\in\cA(A_\eps),\ y_D\in\cA(D),
\end{equation}
on $\cM_\eps:=\cA(A_\eps)\vee\cA(D)$, i.e.\ $\widetilde\omega_\eps=\omega\circ\alpha_s$ with
$\alpha_s:=\mathrm{id}_{\cA(A_\eps)}\otimes\Ad U(s)|_{\cA(D)}$.  Its Lemma~2.1 computes
$\frac{d}{ds}\widetilde\omega_\eps(X)|_{s=0}=i\,\omega([T(g),X])$ with $g=f\chi$.
The task is to exhibit \eqref{eq:N7-state} in the form $\omega\circ\Ad u_s$ so that Theorem~\ref{thm:main} applies.

\begin{citedbox}{Carpi--Weiner 2004, Definition 2.1 (diffeomorphism covariance)}
\emph{Transcribed:} ``A M\"obius covariant net $(A,U)$ is \emph{diffeomorphism covariant} if there is a strongly
continuous projective unitary representation $V$ of $\mathrm{Diff}^+(S^1)$ on $H_A$ such that for all
$\gamma\in\mathrm{Diff}^+(S^1)$ and $I,J\in\mathcal I$: 1.\ $\gamma\in\mathrm{M\ddot ob}\Rightarrow
\Ad(V(\gamma))=\Ad(U(\gamma))$; 2.\ $\gamma|_I=\mathrm{id}_I\Rightarrow\Ad(V(\gamma))|_{A(I)}=\mathrm{id}_{A(I)}$;
3.\ $\gamma(I)=J\Rightarrow V(\gamma)A(I)V(\gamma)^*=A(J)$.''  Their consequences list also records the
Reeh--Schlieder property ($\Om$ cyclic and separating for every $A(I)$), Haag duality $A(I)'=A(I')$, and
Bisognano--Wichmann $U(\Lambda_I(2\pi t))=\Delta_I^{it}$.
\emph{Locator:} S.~Carpi, M.~Weiner, \emph{On the uniqueness of diffeomorphism symmetry in conformal field theory},
CMP 258 (2005) 203--221, arXiv:math/0407190, p.~6 (Def.~2.1) and p.~4 (axioms and consequences).
Local copy \texttt{refs/P1\_CarpiWeiner\_math-0407190.pdf}.
\emph{Screenshot:}\\
\includegraphics[width=\textwidth]{shots/P1_CW_def21.png}\\
\emph{Use here:} property~2 is the only input to Proposition~\ref{prop:F-implement}; property~3 gives
$\Ad V(h_s)(\cA(D))=\cA(h_sD)$.  \emph{Verdict:} directly applicable.
\end{citedbox}

\begin{citedbox}{Fewster--Hollands 2005, Section 5.1 (self-adjointness of the smeared stress tensor)}
\emph{Transcribed:} ``if $\Gamma f=f$ (i.e., $f\in C^\infty_\Gamma(S^1)$) then $\Theta(f)$ is symmetric on $D(L_0)$
and an application of Nelson's commutator theorem (Theorem X.37 in [44]) shows that $\Theta(f)$ is essentially
self-adjoint on any core of $L_0$.  Henceforth we will use $\Theta(f)$ to denote the unique self-adjoint extension.''
\emph{Locator:} C.~J.~Fewster, S.~Hollands, \emph{Quantum energy inequalities in two-dimensional conformal field
theory}, Rev.\ Math.\ Phys.\ 17 (2005) 577, arXiv:math-ph/0412028, p.~26.  Local copy
\texttt{refs/FewsterHollands05\_math-ph-0412028.pdf}.
\emph{Screenshot:}\\
\includegraphics[width=\textwidth]{shots/P1_FH_esa.png}\\
\emph{Use here:} $T(g)$ is self-adjoint and $\Om$, a finite-energy vector ($L_0\Om=0$, $\Om\in D_0$), lies in its
domain; hence $s\mapsto e^{isT(g)}$ is a strongly continuous one-parameter unitary group and Lemma~\ref{lem:tangent}
applies with $A=T(g)$.  \emph{Verdict:} applicable to any diffeomorphism covariant net through its Virasoro subnet;
the normalisation of $\Theta$ differs from Note~7's $T$ by the factor discussed in Note~7 Section~1
(``the BPZ field $2\pi T$''), which does not affect self-adjointness or the domain statement.
\end{citedbox}

\subsection{Unitary implementation}

\begin{proposition}[Implementation of an endomorphism]\label{prop:endo}
Let $\cM$ be a von Neumann algebra with cyclic separating $\Om$ and $\gamma:\cM\to\cM$ a normal unital injective
$*$-endomorphism such that $\Om$ is cyclic for $\gamma(\cM)$.  Then there is a unitary $W$ on $\Hil$ with
$WXW^*=\gamma(X)$ for all $X\in\cM$ and $W^*\Om=\xi$, the natural-cone representative of $\omega\circ\gamma$;
in particular $\omega\circ\gamma=\ip{W^*\Om}{\cdot\,W^*\Om}|_{\cM}$.
\end{proposition}
\lstep{1}{1}{$\omega\circ\gamma$ is a faithful normal state of $\cM$.}
\lby{$\gamma$ normal unital and $\omega$ normal give normality and $(\omega\circ\gamma)(\one)=1$;
$(\omega\circ\gamma)(X^*X)=0\Rightarrow\gamma(X^*X)=0$ ($\omega$ faithful) $\Rightarrow X^*X=0$ ($\gamma$ injective)}
\lstep{1}{2}{\lpick $\xi\in\mathcal P^\natural$ the unique natural-cone representative of $\omega\circ\gamma$;
then $\xi$ is cyclic and separating for $\cM$.}
\lby{existence and uniqueness of natural-cone representatives of normal states: Bratteli--Robinson Thm.~2.5.31;
$\xi$ separating $\Leftrightarrow$ $\omega_\xi$ faithful, which is $\langle1\rangle1$; and for a vector in the natural
cone, separating $\Leftrightarrow$ cyclic (Bratteli--Robinson Prop.~2.5.30)}
\lstep{1}{3}{$W_0:X\xi\mapsto\gamma(X)\Om$ ($X\in\cM$) is a well-defined isometry of $\cM\xi$ onto $\gamma(\cM)\Om$.}
\lby{$\norm{\gamma(X)\Om}^2=\omega(\gamma(X^*X))=(\omega\circ\gamma)(X^*X)=\norm{X\xi}^2$, which gives both
well-definedness ($X\xi=0\Rightarrow\gamma(X)\Om=0$) and isometry}
\lstep{1}{4}{$W:=\overline{W_0}$ is unitary on $\Hil$.}
\lby{$\cM\xi$ is dense by $\langle1\rangle2$ and $\gamma(\cM)\Om$ is dense by hypothesis; an isometry with dense
domain and dense range extends to a unitary}
\lstep{1}{5}{$WXW^*=\gamma(X)$ and $W^*\Om=\xi$.}
\lby{$WXW^*\gamma(Y)\Om=WXY\xi=\gamma(XY)\Om=\gamma(X)\gamma(Y)\Om$ on the dense set $\gamma(\cM)\Om$, and both
sides are bounded; $W\xi=\gamma(\one)\Om=\Om$}
\lqed{1}{6}
\lby{$\langle1\rangle1$--$\langle1\rangle5$}

\begin{corollary}\label{cor:endo-net}
In the setting of Note~7, $\alpha_s$ is a normal unital injective endomorphism of $\cM_\eps$ for $s\ge0$ small, and
$\Om$ is cyclic for $\alpha_s(\cM_\eps)\supseteq\cA(h_sD)$ by the Reeh--Schlieder theorem; hence
$\widetilde\omega_\eps=\ip{W_s^*\Om}{\cdot\,W_s^*\Om}$ for a unitary $W_s$.
\end{corollary}
\lstep{1}{1}{$\alpha_s$ is a normal unital injective $*$-endomorphism of $\cM_\eps$ into itself.}
\lby{$\Ad U(s)$ is a normal $*$-automorphism of $B(\Hil)$ mapping $\cA(D)$ onto $\cA(h_sD)\subseteq\cA(D)$
(covariance, and $h_sD=(0,L/(1+s))\subseteq D$ for $s\ge0$); it acts trivially on $\cA(A_\eps)$ by construction of
$\alpha_s$ as $\mathrm{id}\otimes\Ad U(s)$ on the tensor factorisation
$\cM_\eps\cong\cA(A_\eps)\bar\otimes\cA(D)$ (which holds because $\overline{A_\eps}\cap\overline D=\emptyset$ and the
net is split --- \textbf{this uses the split property}, see Gap~4)}
\lstep{1}{2}{$\Om$ is cyclic for $\alpha_s(\cM_\eps)$.}
\lby{$\alpha_s(\cM_\eps)\supseteq\alpha_s(\cA(D))=\cA(h_sD)$, and $\Om$ is cyclic for $\cA(J)$ for every nonempty
open interval $J$ (Reeh--Schlieder; Carpi--Weiner p.~4, consequence (i))}
\lqed{1}{3}
\lby{$\langle1\rangle1$, $\langle1\rangle2$, Proposition~\ref{prop:endo}}

\begin{proposition}[The implementer is a strongly continuous group, for $\eps>0$]\label{prop:F-implement}
Let $\eps>0$ and let $\chi\in C^\infty(\mathbb R)$ with $\chi=1$ on a neighbourhood $N_1$ of $[0,L]$, $\chi=0$ outside
a compact set $\subseteq(-\eps,\infty)$, and put $g=f\chi\in C^\infty_c(\mathbb R)\subseteq C^\infty(S^1)$.
Let $k_s$ be the flow of the vector field $g\partial_x$ and $u_s:=V(k_s)$.  Then, for $|s|$ small,
\begin{enumerate}[leftmargin=1.7em]
\item $s\mapsto k_s$ is a one-parameter group in $\mathrm{Diff}^+(S^1)$ and $u_s=e^{isT(g)}$ is a strongly continuous
one-parameter unitary group with $T(g)=T(g)^*$ and $\Om\in D(T(g))$;
\item $\Ad u_s|_{\cA(A_\eps)}=\mathrm{id}$ and $\Ad u_s|_{\cA(D)}=\Ad U(s)|_{\cA(D)}$;
\item hence $\Ad u_s|_{\cM_\eps}=\alpha_s$ and $\widetilde\omega_\eps=\ip{u_s^*\Om}{\cdot\,u_s^*\Om}\big|_{\cM_\eps}$,
so hypothesis \textup{(V)} holds with $a=T(g)\Om$ and Theorem~\ref{thm:main} applies:
\begin{equation}\label{eq:F-conclusion}
 -\log\Fid_{\cM_\eps}(\omega,\widetilde\omega_\eps)=\frac{s^2}{8}\gB^{(\eps)}+o(s^2),\qquad
 \gB^{(\eps)}=2\bigl\|(1-J_\eps)(1+\Delta_\eps)^{-1/2}T(g)\Om\bigr\|^2,
\end{equation}
$(\Delta_\eps,J_\eps)$ the modular data of $(\cM_\eps,\Om)$.
\end{enumerate}
\end{proposition}
\lstep{1}{1}{(1).}
\lby{$g\partial_x$ is a smooth vector field on $S^1$ (compactly supported in $\mathbb R$, hence extending by $0$ across
$\infty$), so its flow is a one-parameter group of diffeomorphisms; $V$ is a strongly continuous projective
representation (Carpi--Weiner Def.~2.1), and the generator of $s\mapsto V(k_s)$ in the vacuum representation is
$T(g)$, essentially self-adjoint by Fewster--Hollands \S5.1 with $\Om$ in its domain (finite-energy vector);
a phase choice makes $V(k_s)$ a genuine group because $s\mapsto k_s$ is a one-parameter subgroup and the multiplier
can be trivialised on it (equivalently: define $u_s:=e^{isT(g)}$ by Stone's theorem and note $\Ad u_s=\Ad V(k_s)$)}
\lstep{1}{2}{$\Ad u_s|_{\cA(A_\eps)}=\mathrm{id}$.}
\lby{$k_s|_{A_\eps}=\mathrm{id}$ because $g=0$ on a neighbourhood of $\overline{A_\eps}=(-\infty,-\eps]$
(the support of $\chi$ is a compact subset of $(-\eps,\infty)$), so points of $A_\eps$ are fixed by the flow;
then Carpi--Weiner Def.~2.1(2) with $I\supseteq A_\eps$}
\lstep{1}{3}{$k_s|_{[0,L]}=h_s|_{[0,L]}$ for $|s|$ small.}
\lby{on $N_1\supseteq[0,L]$ one has $g=f$, so the two flows solve the same ODE; $h_s([0,L])=[0,L/(1+s)]\subseteq N_1$
for $|s|$ small; uniqueness of solutions of the ODE $\dot x=f(x)$ with the same initial condition, valid as long as
the trajectory remains in $N_1$}
\lstep{1}{4}{$\Ad u_s|_{\cA(D)}=\Ad U(s)|_{\cA(D)}$.}
\lby{$\gamma_s:=k_s\circ h_s^{-1}$ is the identity on a neighbourhood of $h_s([0,L])$ by $\langle1\rangle3$
(choose $N_1$ so that $k_s=h_s$ on a neighbourhood of $[0,L]$), hence on an interval $I_0\supseteq h_sD$; so
$\Ad V(\gamma_s)|_{\cA(I_0)}=\mathrm{id}$ by Carpi--Weiner Def.~2.1(2), and $\cA(h_sD)\subseteq\cA(I_0)$ by isotony;
finally $\Ad V(k_s)=\Ad V(\gamma_s)\circ\Ad V(h_s)$ (the multiplier cancels in $\Ad$) and
$\Ad V(h_s)(\cA(D))=\cA(h_sD)$}
\lstep{1}{5}{(3).}
\lby{$\langle1\rangle2$, $\langle1\rangle4$ and the fact that $\cM_\eps$ is generated by $\cA(A_\eps)\cup\cA(D)$, so
two normal $*$-homomorphisms agreeing on both generators agree on $\cM_\eps$; then Lemma~\ref{lem:tangent} with
$A=T(g)$ and Theorem~\ref{thm:main}}
\lqed{1}{6}
\lby{$\langle1\rangle1$--$\langle1\rangle5$}

\section{What remains open}\label{sec:open}

\begin{remark}[The upper bound needs no vector representative]\label{rem:upper-no-vector}
Inspection of the proof of Theorem~\ref{thm:upper} shows that it uses \emph{only}: $\omega_s$ normal states and
$\omega_s(y)=\omega(y)+s\varphi(y)+o(s)$ for each $y\in\cM$.  It never uses a vector $a$.  Therefore, defining for an
arbitrary real-linear $\varphi:\Msa\to\mathbb R$ with $\varphi(\one)=0$
\begin{equation}\label{eq:gB-functional}
 \gB(\varphi):=4\sup_{K\in\Msa}\bigl[\varphi(K)-\Var_\omega(K)\bigr]\in[0,+\infty],
\end{equation}
one has, with no further hypothesis,
$\liminf_{s\to0}s^{-2}\bigl(1-\Fid(\omega,\omega_s)\bigr)\ge\tfrac18\gB(\varphi)$
(and hence $\Fid\to0$ faster than any $s^2$-law if $\gB(\varphi)=\infty$).
By Theorem~\ref{thm:three}, $\gB(\varphi)=\gB(a)$ whenever $\varphi$ admits a representative $a\in\Hil$ with
$\operatorname{Im}\ip\Om a=0$.  The lower bound Theorem~\ref{thm:lower}, by contrast, genuinely needs the
\emph{vector} hypothesis (V).  This is exactly the asymmetry that matters at zero collar, where the natural
representative fails to be a vector (Gap~2).
\end{remark}

\begin{proposition}[Stability of $\gB(\varphi)$ under approximation of the tangent]\label{prop:stability}
Let $\varphi_n,\varphi$ be real-linear on $\Msa$ with $\varphi_n(\one)=\varphi(\one)=0$, and put
$G(\psi):=\tfrac14\gB(\psi)=\sup_{K\in\Msa}[\psi(K)-\Var_\omega(K)]$.
\begin{enumerate}[leftmargin=1.7em]
\item \emph{(Lower semicontinuity.)}  If $\varphi_n(K)\to\varphi(K)$ for every $K\in\Msa$, then
 $G(\varphi)\le\liminf_nG(\varphi_n)$.
\item \emph{(Convergence.)}  If in addition $C:=\limsup_nG(\varphi_n)<\infty$ and the convergence is uniform on
 $\{K\in\Msa:\omega(K)=0,\ \norm{K\Om}^2\le C+1\}$, then $G(\varphi_n)\to G(\varphi)$.
\end{enumerate}
\end{proposition}
\lstep{1}{1}{(1).}
\lby{for fixed $K$, $\varphi(K)-\Var_\omega(K)=\lim_n[\varphi_n(K)-\Var_\omega(K)]\le\liminf_nG(\varphi_n)$;
take the supremum over $K$}
\lstep{1}{2}{(2).}
\lproof
\lstep{2}{1}{For $n$ large there is $K_n\in\Msa$ with $\omega(K_n)=0$,
$\varphi_n(K_n)-\norm{K_n\Om}^2\ge G(\varphi_n)-1/n$ and $\norm{K_n\Om}^2\le C+1$.}
\lby{the supremum may be taken over $\omega(K)=0$ (Theorem~\ref{thm:three}, $\langle1\rangle4$--$\langle1\rangle5$);
if $\norm{K\Om}^2>C+1\ge G(\varphi_n)+1$ for a near-maximiser then
$\varphi_n(K)-\norm{K\Om}^2\le G(\varphi_n)$ forces $\varphi_n(K)>C+1$, and replacing $K$ by $tK$, $t\in(0,1)$,
increases $\varphi_n(tK)-t^2\norm{K\Om}^2$ above $G(\varphi_n)$ for suitable $t$ unless
$\norm{K\Om}^2\le G(\varphi_n)+1/n\le C+1$; explicitly, the one-dimensional maximum of
$t\mapsto t\varphi_n(K)-t^2\norm{K\Om}^2$ is $\varphi_n(K)^2/(4\norm{K\Om}^2)\le G(\varphi_n)$, whence
$\varphi_n(K)\le2\norm{K\Om}\sqrt{G(\varphi_n)}$ and
$\varphi_n(K)-\norm{K\Om}^2\ge G(\varphi_n)-1/n$ gives
$\norm{K\Om}^2-2\norm{K\Om}\sqrt{G(\varphi_n)}+G(\varphi_n)\le1/n$, i.e.\
$(\norm{K\Om}-\sqrt{G(\varphi_n)})^2\le1/n$}
\lstep{2}{2}{$G(\varphi)\ge G(\varphi_n)-1/n-\delta_n$ with
$\delta_n:=\sup\{|\varphi(K)-\varphi_n(K)|:\omega(K)=0,\norm{K\Om}^2\le C+1\}\to0$.}
\lby{$G(\varphi)\ge\varphi(K_n)-\norm{K_n\Om}^2=[\varphi_n(K_n)-\norm{K_n\Om}^2]+[\varphi(K_n)-\varphi_n(K_n)]
\ge G(\varphi_n)-1/n-\delta_n$}
\lqed{2}{3}
\lby{$\langle2\rangle2$ gives $G(\varphi)\ge\limsup_nG(\varphi_n)$; with (1), $G(\varphi)=\lim_nG(\varphi_n)$}
\lqed{1}{3}
\lby{$\langle1\rangle1$, $\langle1\rangle2$}

\subsection*{The gaps, precisely}

\begin{verdictgap}
\textbf{Gap 1 (Kubo--Mori beyond finite dimensions).}  Proposition~\ref{prop:KM} proves
$D(\omega\Vert\omega_s)=\frac{s^2}2\gKM+O(s^3)$ only for $\cM=M_n(\mathbb C)$.  For a general von Neumann algebra
Araki's relative entropy $D(\omega\Vert\omega_s)=-\ip{\Om}{\log\Delta_{\omega_s,\omega}\Om}$ has no variational
formula of Alberti type, and the two ingredients that made the fidelity argument work are both missing:
(i) there is no known analogue of $\Fid=\inf_{x>0}\sqrt{\omega(x)\omega_s(x^{-1})}$ giving a \emph{one-sided} bound
from a single $x\in\cM$ with a computable second variation --- the closest is the variational expression
$D(\omega\Vert\varrho)=\sup_{h=h^*\in\cM}\{\omega(h)-\log\varrho(e^h)\}$ (Kosaki's formula / Petz's variational
principle), whose second variation at $h=0$ gives $\sup_h[\varphi(h)-\Var_\omega(h)]$ \emph{again}, i.e.\
$\tfrac14\gB$ and not $\tfrac12\gKM$ --- the golden--Thompson-type gap between the two is exactly the difference
between the Bures and Kubo--Mori metrics; and (ii) monotone convergence along $\cM_n\uparrow\cM$ holds for relative
entropy (Araki 1977, Thm.~3.9), so the analogue of Remark~\ref{rem:hyperfinite} gives
$\liminf_s s^{-2}D\ge\frac12\gKM$ in the hyperfinite case, but the reverse inequality (an upper bound on $D$) is
\emph{not} available; relative entropy is only lower semicontinuous.  \emph{What is missing precisely:} an upper
bound $D(\omega\Vert\omega_s)\le\frac{s^2}2\gKM+o(s^2)$ for a curve $\omega_s=\ip{e^{-isA}\Om}{\cdot\,e^{-isA}\Om}$
with $\Om\in D(A)$ (or $D(A^2)$), for $\cM$ of type III.  Note~7 uses $\gKM$ only through the finite-dimensional
free-fermion computation of Note~5, so Theorem~A of Note~7 is unaffected; Theorem~B's Kubo--Mori half remains
conditional.
\end{verdictgap}

\begin{verdictgap}
\textbf{Gap 2 (the zero-collar, $\xi$-frame representative is not a vector).}  In the modular frame of $I$ the
zero-collar generator is $\widetilde g(\xi)=-4\sinh^2(\xi/2)\mathbf 1_{\xi>0}$ (Note~7, eq.~(11)), and
$\norm{T(g)\Om}^2=\frac1{2\pi}\int\widehat W(k)|\widehat g(k)|^2dk$ \emph{diverges}: $|\widehat g(k)|^2=4/(k^2(k^2+1)^2)
\sim4/k^2$ as $k\to0$ while $\widehat W(0)=\frac{c}{48\pi^2}\ne0$ (verified by agent~V6 of this review).  The two
weighted integrals of Note~7 converge because the Bures weight $(\lambda-1)^2/(\lambda+1)$ and the Kubo--Mori weight
$(\lambda-1)\log\lambda$ both vanish quadratically at $\lambda=1$ ($k=0$).  Consequently:
Theorem~\ref{thm:lower} (lower bound on $\Fid$) does \emph{not} apply directly at zero collar, since hypothesis (V)
requires $a\in\Hil$; Theorem~\ref{thm:upper} does apply, in the form of Remark~\ref{rem:upper-no-vector}, as soon as
the tangent functional exists.  \emph{What is missing precisely:} (a) a proof that $\gB(\varphi)$ of
\eqref{eq:gB-functional} is finite for the zero-collar tangent functional $\varphi$, and (b) its identification with
Note~7's spectral integral $2\int\frac{(\lambda-1)^2}{\lambda+1}d\mu(\lambda)
=\frac1\pi\int\widehat W(k)\frac{(e^{2\pi k}-1)^2}{e^{2\pi k}+1}|\widehat g(k)|^2dk$.  By
Proposition~\ref{prop:stability} it suffices to produce cutoffs $g_n\in C^\infty_c$ (vanishing near the far endpoint
$L$) with $a_n=T(g_n)\Om\in\Hil$ such that (i) $\varphi_n\to\varphi$ pointwise on $\Msa$, (ii)
$\sup_n\gB(a_n)<\infty$, and (iii) the convergence is uniform on Bures-balls
$\{K:\omega(K)=0,\norm{K\Om}\le R\}$; (i) and (ii) look accessible (the spectral integrals converge by dominated
convergence), (iii) is the real work.  A norm-differentiable curve of \emph{vectors} at zero collar would remove the
need for (iii) for the lower bound as well.
\end{verdictgap}

\begin{verdictgap}
\textbf{Gap 3 (the collar limit).}  For $\eps>0$ Proposition~\ref{prop:F-implement} gives the full expansion with
$\gB^{(\eps)}$ computed from the modular data of $(\cM_\eps,\Om)$, which are \emph{not} geometric.  Note~7 evaluates
instead with the geometric modular data of $\cA(I)$, $I=(-\infty,L)$.  The bridge is:
(i) $\cM_\eps$ increases as $\eps\downarrow0$ and $\bigvee_{\eps>0}\cM_\eps=\cA(A)\vee\cA(D)$, which equals
$\cA(I)$ \emph{if the net is strongly additive} (true for the free fermion and for all completely rational nets, but
an extra hypothesis in general); (ii) the tangent functionals $\varphi_\eps$ are the restrictions of one functional
$\varphi$ (Note~7 Lemma~2.1: the representative $g$ may be changed by anything supported away from $D$, and a cutoff
admissible for $\eps$ is admissible for every $\eps'<\eps$), so Proposition~\ref{prop:mono} applies and gives
$\gB^{(\eps)}\uparrow\gB(\varphi)$ with $\gB(\varphi)$ the zero-collar Bures form of \eqref{eq:gB-functional}
on $\cA(I)$.  Thus the collar limit of the \emph{metric} is settled, monotonically; what is not settled is again the
identification of $\gB(\varphi)$ with the spectral integral, i.e.\ Gap~2.  The collar limit of the
\emph{fidelity itself} at fixed $s$ is Note~4 Lemma~10.1.
\end{verdictgap}

\begin{verdictgap}
\textbf{Gap 4 (side hypotheses in (F)).}  Corollary~\ref{cor:endo-net}, $\langle1\rangle1$, writes
$\cM_\eps\cong\cA(A_\eps)\bar\otimes\cA(D)$, which is the \emph{split property} for the pair
$\overline{A_\eps}\cap\overline D=\emptyset$; it holds for conformal nets satisfying the trace-class condition
$\Tr e^{-\beta L_0}<\infty$ (Buchholz--D'Antoni--Longo) but is an extra hypothesis in general.  The split property is
used only to make sense of Note~7's ``$\mathrm{id}\otimes\Ad U(s)$''; Proposition~\ref{prop:F-implement} does not use
it (it constructs $u_s$ directly and shows $\Ad u_s$ is the identity on $\cA(A_\eps)$ and $\Ad U(s)$ on $\cA(D)$,
which \emph{defines} the deformation without any tensor factorisation).  Likewise, the statement that local algebras
of conformal nets are the hyperfinite type~III$_1$ factor (used nowhere in this note, and mentioned in the task) is
standard but its references (Fredenhagen; Gabbiani--Fr\"ohlich; Buchholz--D'Antoni--Fredenhagen) were
\textbf{NOT OBTAINED} here; our proof of Theorem~\ref{thm:main} does not need it.
\end{verdictgap}

\begin{verdictgap}
\textbf{Gap 5 (differentiability of the state curve for other presentations).}  If the deformation is presented only
as a state curve (e.g.\ through natural-cone representatives $\xi_s$ obtained from Proposition~\ref{prop:endo}), then
hypothesis (V) requires \emph{norm} differentiability of $s\mapsto\xi_s$, which does not follow from norm
differentiability of $s\mapsto\omega_s$ in $\cM_*$.  For $\eps>0$ this is bypassed by
Proposition~\ref{prop:F-implement} ($\xi_s=e^{-isT(g)}\Om$ is norm differentiable because $\Om\in D(T(g))$).
In general, the implication ``$\omega_s$ norm-differentiable in $\cM_*$ $\Rightarrow$ the natural-cone
representatives are norm-differentiable'' is \emph{false} without extra assumptions
(the Powers--St\o rmer inequality gives only $\norm{\xi_s-\Om}^2\le\norm{\omega_s-\omega}$, i.e.\ H\"older
exponent $\tfrac12$); the fidelity itself is $1-O(s)$ then, not $1-O(s^2)$.
\end{verdictgap}

\section{Summary}\label{sec:summary}

\begin{center}\small
\begin{tabular}{@{}p{0.42\textwidth}p{0.52\textwidth}@{}}\toprule
\textbf{Statement} & \textbf{Status}\\\midrule
$\varphi(X)=i\ip{A\Om}{X\Om}-i\ip{X^*\Om}{A\Om}$, $\Om\in D(A)$ & Lemma~\ref{lem:tangent}: proved\\
$\eta=i(S^*-1)A\Om$ (Note~7 Lemma~3.1) & Lemma~\ref{lem:eta}: proved \emph{under} $A\Om\in D(\Delta^{-1/2})$, a
hypothesis missing in Note~7\\
$\gB=2\norm{(1-J)(1+\Delta)^{-1/2}A\Om}^2$ & Lemma~\ref{lem:eta}(3): proved, no hypothesis\\
$\tfrac14\gB=\dist(A\Om,\overline{\Mpsa\Om})^2=\sup_{K\in\Msa}[\varphi(K)-\Var_\omega(K)]$ &
Theorem~\ref{thm:three}: proved, general $\cM$, exact\\
$c\in\mathbb R$ vs.\ $c\in\mathbb C$ & Proposition~\ref{prop:c-real}: $c$ redundant over $\mathbb R$; $\mathbb C$ gives
the same value iff $\operatorname{Im}\ip\Om{A\Om}=0$, automatic for $A=A^*$\\
$\Fid=\sup_{v'\in\U(\cM')}|\ip{\Om}{v'u_s^*\Om}|$ & Corollary~\ref{cor:uhlmann}: proved (Alberti--Uhlmann Thm.~2(3))\\
$\Fid\ge1-\frac{s^2}8\gB+o(s^2)$ & Theorem~\ref{thm:lower}: proved, general\\
$\Fid\le1-\frac{s^2}8\gB+o(s^2)$ & Theorem~\ref{thm:upper}: proved, general (Alberti's $\inf_{x>0}$)\\
$-\log\Fid=\frac{s^2}8\gB+o(s^2)$ (Note~7 Lemma~3.2, fidelity) & Theorem~\ref{thm:main}: \textbf{proved}\\
$\gB^{(n)}\uparrow\gB$ along $\cM_n\uparrow\cM$ & Proposition~\ref{prop:mono}: proved\\
$D(\omega\Vert\omega_s)=\frac{s^2}2\gKM+o(s^2)$ & Proposition~\ref{prop:KM}: proved in finite dimensions only
(Gap~1)\\
Note~7's curve is $\omega\circ\Ad u_s$, $u_s$ a strongly continuous group & Proposition~\ref{prop:F-implement}:
proved for $\eps>0$, $u_s=e^{isT(g)}$\\
Zero collar ($T(g)\Om\notin\Hil$) & Gaps~2,~3\\\bottomrule
\end{tabular}
\end{center}

The proof of Theorem~\ref{thm:main} may be summarised in one line.  The Bures distance is
$\inf_{v'\in\U(\cM')}\norm{\Om-v'\xi_s}$, whose square is $s^2$ times the squared distance from the tangent
$A\Om$ to the \emph{gauge directions} $\Mpsa\Om$; that distance is computed exactly by the bounded operator
$(1+\Delta)^{-1}(1-\Delta^{1/2}J)$; and it is dual, in the sense of the elementary Legendre identity
$\sup_k[2\operatorname{Re}\ip wk-\norm k^2]=\norm{P w}^2$, to a supremum over $\Msa$ that pairs with Alberti's
inequality $\Fid^2\le\omega(x)\,\omega_s(x^{-1})$.  The two sides of Alberti's variational formula therefore give the
two sides of the expansion, and the modular formula of Note~7 sits exactly in the middle.

\begin{thebibliography}{99}
\bibitem{AU02} P.~M.~Alberti, A.~Uhlmann, \emph{On Bures distance and $*$-algebraic transition probability between
inner derived positive linear forms}, arXiv:math-ph/0202038 (2002); Rep.\ Math.\ Phys.\ \textbf{46} (2000) 97.
Local copy \texttt{refs/AlbertiUhlmann02\_math-ph-0202038.pdf}.  Theorem~2 there restates
P.~M.~Alberti, \emph{A note on the transition probability over C$^*$-algebras}, Lett.\ Math.\ Phys.\ \textbf{7}
(1983) 25 (\textbf{NOT OBTAINED}, paywalled) and A.~Uhlmann, \emph{The ``transition probability'' in the state space
of a $*$-algebra}, Rep.\ Math.\ Phys.\ \textbf{9} (1976) 273 (\textbf{NOT OBTAINED}, paywalled).
\bibitem{BR1} O.~Bratteli, D.~W.~Robinson, \emph{Operator Algebras and Quantum Statistical Mechanics 1}, 2nd ed.,
Springer 1987: Prop.~2.5.3, Prop.~2.5.9, Thm.~2.4.16 (Kaplansky), Thm.~2.5.14 (Tomita), Prop.~2.5.30,
Thm.~2.5.31 (natural cone).  \textbf{NOT OBTAINED} as a file; the statements used are standard and are quoted
individually at the point of use.
\bibitem{CW} S.~Carpi, M.~Weiner, \emph{On the uniqueness of diffeomorphism symmetry in conformal field theory},
Commun.\ Math.\ Phys.\ \textbf{258} (2005) 203, arXiv:math/0407190.
Local copy \texttt{refs/P1\_CarpiWeiner\_math-0407190.pdf}.
\bibitem{FH05b} C.~J.~Fewster, S.~Hollands, \emph{Quantum energy inequalities in two-dimensional conformal field
theory}, Rev.\ Math.\ Phys.\ \textbf{17} (2005) 577, arXiv:math-ph/0412028.
Local copy \texttt{refs/FewsterHollands05\_math-ph-0412028.pdf}.
\bibitem{Araki77b} H.~Araki, \emph{Relative entropy for states of von Neumann algebras II}, Publ.\ RIMS \textbf{13}
(1977) 173: Thm.~3.9 (martingale convergence).  Local copy \texttt{refs/Araki77\_PRIMS13.pdf}.
\bibitem{RS1} M.~Reed, B.~Simon, \emph{Methods of Modern Mathematical Physics I}, Thm.~VIII.7 (Stone).
\textbf{NOT OBTAINED} as a file; standard.
\bibitem{N4b} Note~4, \texttt{continuum\_petz\_free\_fermion.tex} (Lemma~10.1).
\bibitem{N5b} Note~5, \texttt{fidelity\_recovered\_quasifree.tex} (Prop.~2.3, Thm.~3.1).
\bibitem{N7b} Note~7, \texttt{universality\_normalization.tex} (Lemmas~2.1, 3.1, 3.2; \S8).
\end{thebibliography}

\end{document}
